\chapter{Strategy observables and the expanded space}
\label{chap:qpbt-observables}
% This chapter covers the first three subsections of the source appendix
% ``Analysis of the Pauli basis test''
% (references/qpbt-paper/14_analysis_of_the_pauli_basis_test.tex, lines 1-679).
% The paper's subsection labels sec:qld-prelim, sec:commutation and
% sec:expanding are kept on the sections below for cross-reference fidelity.
% The source restates thm:pauli verbatim at the head of the appendix under the
% label thm:pauli-appendix (lines 6-20); the statement is the same theorem, and
% all references here are to thm:pauli as stated in Chapter chap:qpbt-test.

This chapter and the two that follow contain the proof of the soundness
theorem for the Pauli basis test, Theorem~\ref{thm:pauli}, stated in
Chapter~\ref{chap:qpbt-test}. Starting from a strategy that succeeds in the
game $\mathfrak{G}^{\mathrm{Pauli}}_{(q,m,d)}$ with high probability, the proof
constructs a representation of the Pauli group on the players' Hilbert space
in several stages. In Section~\ref{sec:commutation} the strategy measurements
are used to define a set of observables that correspond to a subset of the
Pauli $X$ and $Z$ observables, and these observables are shown to
approximately satisfy the associated Pauli group relations. In
Section~\ref{sec:expanding}, ancillary registers containing EPR states are
adjoined to the players' state, and a new set of approximately commuting
observables acting on the expanded state is defined.
Chapter~\ref{chap:qpbt-combining} combines the approximately commuting
observables into a strategy for the \emph{classical} low individual degree
test and applies the soundness theorem of that test to construct a single
projective measurement that simultaneously measures all of them.
Chapter~\ref{chap:qpbt-extraction} uses this simultaneous measurement to
construct an exact representation of the Pauli group, shows that it is close
to the original strategy observables on the expanded state, and deduces
Theorem~\ref{thm:pauli}. Section~\ref{sec:qld-prelim} below collects
preliminary definitions and lemmas used throughout the analysis.

\section{Preliminaries}
\label{sec:qld-prelim}

\begin{definition}[Individual and total degree polynomial classes]
	\label{def:ideg-deg-polynomials}
	\lean{MIPStarRE.LDT.Preliminaries.polyFunc}
	\lean{MIPStarRE.QPBT.TotalDegreePoly}
	\lean{MIPStarRE.QPBT.DegPoly, MIPStarRE.QPBT.evalCoefficient}
	\lean{MIPStarRE.QPBT.DegPoly.padTo, MIPStarRE.QPBT.EvaluatesTo, MIPStarRE.QPBT.evalOpt}
	\leanok
	\uses{def:line, def:line-representative, def:polyfunc}
	Given integers $d, m \geq 0$ and a prime power $q$, denote by
	$\mathrm{ideg}_{d,m}(\F_q)$ the set of polynomials in $m$ variables over
	$\F_q$ with \emph{individual} degree at most $d$ in each variable. When
	the field is clear from context we write $\mathrm{ideg}_{d,m}$, and when
	the number of variables is also clear from context, $\mathrm{ideg}_d$. We
	denote by $\deg_d$ the set of polynomials with \emph{total} degree at
	most $d$.
	% The line-restricted classes below are introduced only implicitly in the
	% source, in the statement of the lemma labelled lem:qld-comm-line-cons,
	% which reads them as polynomials "defined on the line". The coefficient
	% reading adopted here follows the answer format of the decision procedure
	% of the low individual degree game (source figure fig:ld-decider) and the
	% representative convention of def:polyfunc; see the remark following this
	% definition.
	For a line $\ell \subseteq \F_q^m$ (Definition~\ref{def:line}) presented
	by a base point $u_0$ and a direction $v$ --- for the lines arising as
	question contents, the pair read off the line's description, whose base
	point is the canonical representative of
	Definition~\ref{def:line-representative} --- and an integer $c \geq 0$,
	we write $\deg_c(\ell)$ for the set of univariate polynomials of degree
	at most $c$, each given by its list of $c+1$ coefficients and identified
	by its polynomial representative rather than by the function it induces,
	as in Definition~\ref{def:polyfunc}; the argument $\ell$ records the
	line to which such an answer refers, through the parameterization
	$t \mapsto u_0 + tv$. The cases used below are $\deg_d(\ell)$ and
	$\deg_{md}(\ell)$; whenever $d \leq md$, extending coefficient lists by
	zeros gives the inclusion $\deg_d(\ell) \subseteq \deg_{md}(\ell)$. For
	$f \in \deg_c(\ell)$, $u \in \ell$, and $a \in \F_q$, we say that $f$
	\emph{evaluates to $a$ at $u$}, written $f(u) = a$, if $f(t) = a$ for
	every $t \in \F_q$ with $u = u_0 + tv$. When $v \neq 0$ this $t$ is
	unique, so every $f$ evaluates at every point of $\ell$; when $v = 0$,
	so that $\ell = \{u_0\}$, exactly those $f$ that induce a constant
	function on $\F_q$ evaluate at $u_0$, to their constant value. An
	operator indexed by $f(u)$ is understood to be the zero operator when
	$f$ does not evaluate at $u$, and for a family $\{N_f\}$ indexed by
	$f \in \deg_c(\ell)$ we write
	$N_{[\mathrm{eval}_u(\cdot)=a]} = \sum_{f \colon f(u) = a} N_f$; when
	$v = 0$ the classes indexed by $a \in \F_q$ form a sub-measurement
	rather than a partition of the outcomes. The \emph{completed} family of
	evaluation classes is the family indexed by $a \in \F_q \cup \{\bot\}$
	obtained by adjoining the additional class
	$N_{[\mathrm{eval}_u(\cdot)=\bot]} = \sum_{f} N_f$, the sum over those
	$f$ that do not evaluate at $u$; when $\{N_f\}$ is a POVM,
	respectively a projective measurement, the completed classes form one
	as well, and when $v \neq 0$ the class
	$N_{[\mathrm{eval}_u(\cdot)=\bot]}$ is the zero operator.
\end{definition}

\begin{remark}[Line answers as coefficient lists]
	\label{rem:deg-line-representatives}
	The source refers to the elements of $\deg_{md}(\ell)$ as polynomials
	``defined on the line $\ell$'', that is, as functions on the point set
	of $\ell$, while its decision procedure expects each line answer as a
	list of coefficients. The two readings differ. When $md \geq q$,
	distinct coefficient lists induce the same function $\F_q \to \F_q$;
	and when $\ell$ has direction $v = 0$ --- a case the question
	distributions produce with positive probability
	(Remark~\ref{rem:ld-win-zero-direction}) --- the point set of $\ell$ is
	the singleton $\{u_0\}$, so a function on $\ell$ retains only a single
	field element of the answer. Definition~\ref{def:ideg-deg-polynomials}
	therefore adopts the coefficient reading, under which the strategy's
	line measurements are indexed by exactly the answers of the prescribed
	form, and its evaluation convention reproduces the universal
	quantification of the line-versus-point clauses of the win predicates
	(Definition~\ref{def:ld-win-predicate}).
\end{remark}

Throughout this chapter $q = 2^k$ with $k$ odd is an admissible field size
(Definition~\ref{def:admissible-size}) and $M = 2^m$ is the number of
$q$-ary qudits tested for per player, as in
Chapters~\ref{chap:qpbt-test} and~\ref{chap:qpbt-extraction}. (The letter
$M$ also denotes the measurement operators of the strategy under analysis,
as in the tuple $\mathcal{S} = (\ket{\psi}, M)$ below and in the
super- and subscripted individual operators such as
$M^{(\mathrm{Point},W),u}_a$; the meaning is determined by context, as in
the source.) Recall from
Chapter~\ref{chap:qpbt-algebra} the low-degree code
(see Definition~\ref{def:low-degree-encoding}), which maps each vector
$h \in \F_q^{M}$, with
coordinates indexed by binary strings of length $m$, to a polynomial
$g_h \in \mathrm{ideg}_{d,m}(\F_q)$ such that $g_h(u) = h_u$ for every
$u \in \{0,1\}^m$; for $x \in \F_q^m$ the associated vector
$\mathrm{ind}_m(x) \in \F_q^{M}$ satisfies
$g_h(x) = h \cdot \mathrm{ind}_m(x)$ for all $h$. Recall also the finite
field trace $\tr = \tr_{q \to 2} \colon \F_q \to \F_2$
(see Definition~\ref{def:subfield-trace}).

\begin{definition}[Commuting and anticommuting tuples]
	\label{def:anticommuting-tuple}
	\lean{MIPStarRE.QPBT.PauliTuple}
	\lean{MIPStarRE.QPBT.IsAnticommuting, MIPStarRE.QPBT.IsCommuting}
	\leanok
	% The quantity gamma is introduced in the source at eq:gamma-value, inside
	% the completeness analysis of the Pauli basis test; the definition below
	% is self-contained and does not depend on that occurrence.
	\uses{def:low-degree-encoding, def:subfield-trace}
	For a tuple
	$\omega = (u_X, u_Z, r_X, r_Z) \in (\F_q^m)^2 \times (\F_q)^2$, let
	\[
		\gamma(\omega) \,=\, \tr\big( (r_X \cdot \mathrm{ind}_m(u_X)) \cdot (r_Z \cdot \mathrm{ind}_m(u_Z)) \big)\,,
	\]
	where $r_W \cdot \mathrm{ind}_m(u_W)$ is the scalar multiple of the
	vector $\mathrm{ind}_m(u_W) \in \F_q^{M}$, the middle $\cdot$ is the
	$\F_q$-inner product of two vectors in $\F_q^{M}$, and
	$\tr = \tr_{q \to 2}$. The tuple $\omega$ is called \emph{anticommuting}
	if $\gamma(\omega) \neq 0$, and \emph{commuting} if $\gamma(\omega) = 0$.
\end{definition}

The sign $(-1)^{\gamma(\omega)}$ is exactly the sign with which the
generalized Pauli observables $\tau^X(r_X \cdot \mathrm{ind}_m(u_X))$ and
$\tau^Z(r_Z \cdot \mathrm{ind}_m(u_Z))$ commute
(see Lemma~\ref{lem:twisted-commutation});
this is the source of the terminology.

\begin{lemma}[Probability of anticommuting tuples]
	\label{fact:omega-anticomm-prob}
	\lean{MIPStarRE.QPBT.anticommProb, MIPStarRE.QPBT.anticommProb_eq}
	\lean{MIPStarRE.QPBT.commProb, MIPStarRE.QPBT.commProb_ge_half}
	\lean{MIPStarRE.QPBT.commProb_eq_one_sub_anticommProb}
	\lean{MIPStarRE.QPBT.anticommProb_ge_of_m_le_q}
	\lean{MIPStarRE.QPBT.anticommProb_ge_of_one_le_md}
	\leanok
	% The source states this result in a `fact' environment; it is stated here
	% as a lemma, and the paper label is kept verbatim.
	% The exact value of the anticommuting probability and the constant lower
	% bounds under the hypothesis m <= q are additions relative to the source,
	% which asserts only the two bounds displayed at the end of the statement;
	% the source states those bounds without the hypothesis m, d >= 1, under
	% which alone they hold (see the remark following this lemma). The
	% unconditional constant bounds are what the proof of the lemma labelled
	% lem:qld-win-implications uses.
	\uses{def:anticommuting-tuple}
	A tuple $\omega = (u_X,u_Z,r_X,r_Z)$ sampled uniformly at random from
	$(\F_q^m)^2 \times (\F_q)^2$ is anticommuting with probability exactly
	\[
		\big(1 - q^{-1}\big)^{m+1} \cdot \frac{1}{2}\,,
	\]
	and commuting with the complementary probability, which is at least
	$\frac12$. In particular, if $m \leq q$, then $\omega$ is anticommuting
	with probability at least $2^{-4}$ and commuting with probability at
	least $\frac12$. If $m, d \geq 1$, then the probability that $\omega$ is
	anticommuting is at least
	\[
		(1 - q^{-m}) \cdot (1 - q^{-1}) \cdot \Big(1 - \frac{md}{q}\Big) \cdot \frac{1}{2}
		\,\geq\, \Big(1 - \frac{3md}{q}\Big) \cdot \frac{1}{2}\,,
	\]
	and the probability that $\omega$ is commuting is also at least
	$\big(1 - \frac{3md}{q}\big) \cdot \frac{1}{2}$.  These corrections are
	documented in \cite{gap:qpbt_anticommuting-probability}.
\end{lemma}
\begin{proof}
	\leanok
	\uses{def:low-degree-encoding, def:indicator-vector, def:admissible-size, def:binary-representation, lem:one}
	By the definition of the indicator vector
	(Definition~\ref{def:indicator-vector}), expanding the product over the
	coordinates $i \in \{1, \dots, m\}$ into a sum over $y \in \{0,1\}^m$
	gives
	\[
		\mathrm{ind}_m(u_X) \cdot \mathrm{ind}_m(u_Z)
		\,=\, \sum_{y \in \{0,1\}^m} \mathrm{ind}_{m,y}(u_X)\, \mathrm{ind}_{m,y}(u_Z)
		\,=\, \prod_{i=1}^{m} \big( u_{X,i}\, u_{Z,i} + (1 - u_{X,i})(1 - u_{Z,i}) \big)\,,
	\]
	and since $\F_q$ has characteristic $2$, the $i$-th factor equals
	$1 + u_{X,i} + u_{Z,i}$. Hence
	\[
		\gamma(\omega) \,=\, \tr\Big( (r_X r_Z) \cdot \prod_{i=1}^{m} (1 + u_{X,i} + u_{Z,i}) \Big)\,.
	\]
	For any fixed $u_X$, the $m$ factors $1 + u_{X,i} + u_{Z,i}$ are
	independent and uniformly distributed in $\F_q$ over the choice of
	$u_Z$, so all of them are nonzero with probability $(1-q^{-1})^m$;
	independently, $r_X \neq 0$ with probability $1 - q^{-1}$. If $r_X = 0$
	or some factor vanishes, then $\gamma(\omega) = 0$. Otherwise, over the
	uniform choice of $r_Z$, the argument of the trace is a uniformly random
	element of $\F_q$, and a uniformly random element of $\F_q$ has trace
	$0$ with probability exactly $\frac12$: the first vector $e_1$ of the
	fixed self-dual normal basis satisfies $\tr(e_1) = 1$ by
	Lemma~\ref{lem:one}, the trace
	(Definition~\ref{def:binary-representation}) is additive, and $\F_q$
	has characteristic $2$, so the translation $c \mapsto c + e_1$ is an
	involution of $\F_q$ carrying $\{c : \tr(c) = 0\}$ bijectively onto
	$\{c : \tr(c) \neq 0\}$; the two sets therefore have the same
	cardinality. Combining, the probability that
	$\gamma(\omega) \neq 0$ is exactly $(1-q^{-1})^{m+1} \cdot \frac12$, and
	the probability that $\omega$ is commuting is the complement, which is
	at least $\frac12$.

	If $m \leq q$, then $(1-q^{-1})^{m+1} \geq (1-q^{-1})^{q+1}$, and this
	last quantity is at least $2^{-3}$ in each of the two cases permitted
	by admissibility (Definition~\ref{def:admissible-size}): the exponent
	$k$ with $q = 2^k$ is odd, so either $q = 2$ or $8 \mid q$. If $q = 2$,
	then $m \leq 2$ and $(1-q^{-1})^{m+1} = 2^{-(m+1)} \geq 2^{-3}$. If
	$q = 8j$ with $j \geq 1$, then Bernoulli's inequality gives
	$(1-q^{-1})^{j} \geq 1 - j q^{-1} = \frac78$, while
	$m + 1 \leq q + 1 \leq 9j$, so
	$(1-q^{-1})^{m+1} \geq \big( (1-q^{-1})^{j} \big)^{9} \geq \big( \frac78 \big)^{9} \geq \frac18$.
	Hence the anticommuting probability is at least $2^{-4}$.

	Finally, suppose $m, d \geq 1$, and consider the first inequality of
	the display, which bounds the anticommuting probability from below.
	If $md > q$, that lower bound is negative --- it is the product of
	the positive factors $1 - q^{-m}$, $1 - q^{-1}$, and $\frac12$ with
	the negative factor $1 - md/q$ --- so the anticommuting probability,
	being nonnegative, exceeds it. If $md \leq q$, then by Bernoulli's
	inequality
	$(1-q^{-1})^{m} \geq 1 - m/q \geq 1 - md/q \geq (1-q^{-m})(1-md/q)$,
	where the last step multiplies the nonnegative quantity $1 - md/q$
	by $1 - q^{-m} \leq 1$; multiplying by $1 - q^{-1}$ shows that the
	anticommuting probability is at least the first displayed bound. For
	the second inequality of the display, expanding the product shows
	that $(1-q^{-m})(1-q^{-1})(1-md/q)$ exceeds
	$1 - q^{-m} - q^{-1} - md/q$ by
	$q^{-m-1} + \frac{md}{q} \big( q^{-m} + q^{-1} - q^{-m-1} \big)$,
	a sum of nonnegative terms whatever the sign of $1 - md/q$; since
	$md \geq 1$ gives $q^{-m} \leq q^{-1} \leq md/q$, the product is
	therefore at least $1 - 3md/q$.
	The commuting probability is at least
	$\frac12 \geq \big(1 - \frac{3md}{q}\big) \cdot \frac12$.
\end{proof}

\begin{remark}[Source correction: the anticommuting probability]
	\label{rem:omega-anticomm-prob-correction}
	In the source, the statement corresponding to
	Lemma~\ref{fact:omega-anticomm-prob} asserts only the two bounds of the
	final display, without the hypothesis $m, d \geq 1$, and its proof
	begins with the claims that $\mathrm{ind}_m(u_X)$ vanishes if and only
	if $u_X = 0$ and that $u_X = 0$ forces $\omega$ to be commuting. Both
	claims are incorrect: the coordinates of $\mathrm{ind}_m(u)$ sum to $1$
	for every $u \in \F_q^m$, so $\mathrm{ind}_m(u)$ never vanishes, and for
	$(q,m,d) = (2,1,0)$ the tuple $\omega = (0,0,1,1)$ is anticommuting.
	For $d = 0$ the bounds themselves fail: with $(q,m,d) = (2,1,0)$ ---
	a tuple with admissible field size and $m \mid q$, excluded from
	Definition~\ref{def:admissible} only by the requirement $d \geq 1$
	that the source leaves implicit --- the exact anticommuting
	probability
	$(1-q^{-1})^{m+1} \cdot \frac12 = \frac18$ is smaller than the claimed
	lower bound $\big(1 - \frac{3md}{q}\big) \cdot \frac12 = \frac12$. The
	proof given above therefore computes the probability exactly, using the
	characteristic-two product formula for
	$\mathrm{ind}_m(u_X) \cdot \mathrm{ind}_m(u_Z)$, in place of the
	source's argument via the Schwartz--Zippel lemma; the exact value also
	supplies the unconditional constant lower bounds used in the proof of
	Lemma~\ref{lem:qld-win-implications}.  The full comparison is in
	\cite{gap:qpbt_anticommuting-probability}.
\end{remark}

\begin{definition}[Closeness of question-indexed operators]
	\label{def:approx-question-indexed-operators}
	\lean{MIPStarRE.QPBT.opDistSq}
	\leanok
	\uses{def:povm-distance}
	Let $\{A^x\}$ and $\{B^x\}$ be families of operators (not necessarily
	POVM elements) indexed by questions $x \in \mathcal{X}$ but not by
	answers, let $\mu$ be a distribution over $\mathcal{X}$, and let
	$\ket{\psi}$ be a state, both fixed by context. Write
	$A^x \approx_\delta B^x$ to denote
	\[
		\E_{x \sim \mu} \bra{\psi} (A^x - B^x)^\dagger (A^x - B^x) \ket{\psi} \,\leq\, O(\delta)\,.
	\]
	This extends the answer-indexed closeness of
	Definition~\ref{def:povm-distance} to answer-free families of operators,
	such as observables; as there, the right-hand side allows a hidden
	universal constant.
\end{definition}

The following two lemmas relate the closeness between sets of operators and
weighted sums of those operators.

\begin{lemma}
	\label{lem:avg-closeness}
	\lean{MIPStarRE.QPBT.avg_closeness}
	\leanok
	\uses{def:approx-question-indexed-operators}
	Let $\{A^x\}$ and $\{B^x\}$ be operators indexed by a finite set
	$\mathcal{X}$, let $\mu$ be a probability distribution over
	$\mathcal{X}$, and let $\{\alpha_x\}$ be complex numbers with
	$|\alpha_x| \leq 1$. If $A^x \approx_\delta B^x$ on average over
	$x \sim \mu$, then $A \approx_\delta B$ for the single operators
	$A = \E_{x} \alpha_x A^x$ and $B = \E_{x} \alpha_x B^x$, regarded as
	indexed by a one-element question set; quantitatively,
	% The closeness of Definition def:approx-question-indexed-operators
	% carries a hidden universal constant on both sides; the quantitative
	% sentence records the constant-free inequality the proof establishes.
	\[
		\big\| (A - B) \ket{\psi} \big\|^2
		\,\leq\, \E_x \big\| (A^x - B^x) \ket{\psi} \big\|^2\,,
	\]
	so the conclusion holds with the same implicit constant as the
	hypothesis.
\end{lemma}
\begin{proof}
	\leanok \uses{lem:distance-averaged-square-root}
	Expand
	\[
		\Big\| \E_x \alpha_x (A^x - B^x) \ket{\psi} \Big\|^2
		\,\leq\, \Big( \E_x \big\| (A^x - B^x) \ket{\psi} \big\| \Big)^2
		\,\leq\, \E_x \big\| (A^x - B^x) \ket{\psi} \big\|^2\,,
	\]
	using the triangle inequality followed by Jensen's inequality; the
	right-hand side is at most $O(\delta)$ by the hypothesis, with the same
	implicit constant.
\end{proof}

\begin{lemma}[Formalization support for Lemma~\ref{lem:povm-to-obs}]
	\label{lem:povm-to-obs-operator-families}
	\lean{MIPStarRE.QPBT.povm_to_obs}
	\leanok
	\uses{def:approx-question-indexed-operators, def:povm-distance}
	This formalization-only auxiliary is not a named statement of the paper; it
	isolates the finite-sum estimate used in Lemma~\ref{lem:povm-to-obs}. Let
	$\{A^x_a\}$ and $\{B^x_a\}$ be arbitrary operator families and let
	$|\alpha_a|=1$. Then
	\[
		\E_x \Big\|\sum_a \alpha_a(A^x_a-B^x_a)\ket\psi\Big\|^2
		\leq |\mathcal A|\,\E_x\sum_a
		\big\|(A^x_a-B^x_a)\ket\psi\big\|^2.
	\]
\end{lemma}
\begin{proof}
	\leanok
	Apply the triangle inequality to the sum over $a$, square, and use the
	finite Cauchy--Schwarz inequality. Averaging preserves the resulting
	pointwise inequality.
\end{proof}

\begin{lemma}
	\label{lem:povm-to-obs}
	\lean{MIPStarRE.QPBT.povm_to_obs_of_measurements}
	\leanok
	\uses{def:approx-question-indexed-operators, def:povm-distance}
	Let $\mathcal{A}$ be a finite set, let $\{A^x_a\}_{a \in \mathcal{A}}$
	and $\{B^x_a\}_{a \in \mathcal{A}}$ be POVMs, and let
	$\{\alpha_a\}_{a \in \mathcal{A}}$ be complex numbers on the unit
	circle. Let $A^x = \sum_a \alpha_a A^x_a$ and
	$B^x = \sum_a \alpha_a B^x_a$. Then
	\[
		A^x_a \approx_\delta B^x_a \quad\text{implies}\quad A^x \approx_{\delta'} B^x
	\]
	for $\delta' = |\mathcal{A}| \cdot \delta$, where the hypothesis is the
	answer-indexed closeness of Definition~\ref{def:povm-distance} and the
	conclusion is the question-indexed closeness of
	Definition~\ref{def:approx-question-indexed-operators}.
\end{lemma}
\begin{proof}
	\leanok
	\uses{lem:povm-to-obs-operator-families}
	Apply Lemma~\ref{lem:povm-to-obs-operator-families} to the two families of
	POVM effects.
\end{proof}

The following lemma shows that POVMs that are approximately projective, in
the sense of being self-consistent, are close to exact projective
measurements.

\begin{lemma}[Orthonormalization lemma]
	\label{lem:ortho}
	\lean{MIPStarRE.QPBT.exists_projective_close_of_consistent}
	\leanok
	\uses{def:consistency, def:povm-distance}
	Let $0 \leq \delta \leq 1$. There exists a function
	$\eta_{\mathrm{ortho}}(\delta) = O\big(\delta^{1/4}\big)$ such that the
	following holds. Let $\ket{\psi}$ be a state on
	$\hilb_{\mathsf{A}} \ot \hilb_{\mathsf{B}}$ where $\hilb_{\mathsf{A}}$
	and $\hilb_{\mathsf{B}}$ are finite dimensional, and let $\{Q_a\}$ and
	$\{R_a\}$ be POVMs on $\hilb_{\mathsf{A}}$ and $\hilb_{\mathsf{B}}$
	respectively such that
	\begin{equation}\label{eq:cons-cond}
		Q_a \ot \Id_{\mathsf{B}} \,\simeq_\delta\, \Id_{\mathsf{A}} \ot R_a\,.
	\end{equation}
	Then there exists a projective measurement $\{P_a\}$ on
	$\hilb_{\mathsf{A}}$ such that, with respect to $\ket{\psi}$,
	\begin{equation}\label{eq:cons-ccl}
		P_a \ot \Id_{\mathsf{B}} \,\approx_{\eta_{\mathrm{ortho}}(\delta)}\, Q_a \ot \Id_{\mathsf{B}}\,.
	\end{equation}
\end{lemma}
\begin{proof}
	\leanok
	\uses{lem:ortho-explicit-constant}
	The source imports this lemma: it first appears in~\cite{KV11}, and a
	self-contained proof is given in the reference [ML20] cited
	in~\cite{Ji2020MIPStar}. Here it is obtained from
	Lemma~\ref{lem:ortho-explicit-constant}, which exhibits
	$\eta_{\mathrm{ortho}}(\delta) = 220\,\delta^{1/4}$: passing to Alice's
	reduced state, the rank-reduction and orthonormalization construction of
	Chapter~\ref{chap:projective} produces a projective family on
	$\hilb_{\mathsf{A}}$ close to $\{Q_a\}$, and completing that family at one
	fixed outcome yields the projective measurement $\{P_a\}$.
\end{proof}

\begin{lemma}[Formalization support for Lemma~\ref{lem:ortho}]
	\label{lem:ortho-explicit-constant}
	\lean{MIPStarRE.QPBT.projective_rounding_with_explicit_constant}
	\leanok
	\uses{def:consistency, def:povm-distance}
	This formalization-only auxiliary is not a named statement of the paper; it is
	the explicit-constant form of Lemma~\ref{lem:ortho} from which that lemma is
	obtained. Let $\ket{\psi}$ be a unit vector on
	$\hilb_{\mathsf{A}} \ot \hilb_{\mathsf{B}}$ with $\hilb_{\mathsf{A}}$ and
	$\hilb_{\mathsf{B}}$ finite dimensional, let $\{Q_a\}$ and $\{R_a\}$ be POVMs
	on $\hilb_{\mathsf{A}}$ and $\hilb_{\mathsf{B}}$ respectively, and let
	$\delta \geq 0$ satisfy the exact consistency-defect bound
	\[
		\sum_{\substack{a,b \\ a \ne b}}
		\bra{\psi}
		Q_a \ot R_b \ket{\psi}
		\leq \delta\,.
	\]
	Then there is a projective measurement $\{P_a\}$ on $\hilb_{\mathsf{A}}$ with
	\[
		\sum_{a}
		\left\|\left((P_a-Q_a) \ot \Id_{\mathsf{B}}\right)\ket{\psi}\right\|^2
		\leq 220\,\delta^{1/4}\,.
	\]
\end{lemma}
\begin{proof}
	\leanok
	\uses{lem:projective-low-rank-sum, lem:Q-completeness, def:projective-P,
		lem:sqrt-Q-completeness, lem:P-projectivity, lem:P-Q-approx,
		prop:triangle-inequality-for-approx_delta, def:left-marginal-state,
		lem:matrix-rank-zero, lem:left-marginal-state-evaluation, lem:left-marginal-state-normalization}
	Pass to Alice's reduced state $\rho_{\mathsf A}$. Consistency gives
	$\sum_a \operatorname{Tr}_{\mathrm{norm}}\!\left(
	\rho_{\mathsf A}(Q_a-Q_a^2)\right) \leq 2\delta$.
	When $\delta^{1/4} \leq 1/24$, Lemma~\ref{lem:projective-low-rank-sum}
	replaces $\{Q_a\}$ by projectors. Lemma~\ref{lem:Q-completeness} rules out
	the alternative in which all ranks vanish, so the rank-sum space used by the
	$Q/X/\widehat X/P$ construction is nonzero. Lemma~\ref{lem:sqrt-Q-completeness}
	controls the square root of their total operator on $\rho_{\mathsf A}$.
	Lemma~\ref{lem:P-projectivity} supplies $P$ and its total-operator identity,
	while Lemma~\ref{lem:P-Q-approx} makes $P$ close to $Q$. Completing $P$ at one
	fixed outcome and applying the squared-distance triangle bound gives
	$220\,\delta^{1/4}$, with residual mass controlled by square-root completeness.
	When $\delta^{1/4} > 1/24$, the target exceeds the universal bound $4$, so
	completing the zero sub-measurement already proves the conclusion.
\end{proof}

\section{Strategies}
\label{sec:commutation}

% In the source this setup --- fixing the parameter tuple, the strategy and
% its success probability, and making the measurements projective via
% Naimark's theorem --- is untitled prose at the head of the subsection
% (references/qpbt-paper/14_analysis_of_the_pauli_basis_test.tex, lines
% 155-171), not a numbered environment. It is recorded here as a standing
% convention paragraph together with a lemma for the one mathematical step
% it contains, the reduction to a projective strategy.
\paragraph{The strategy under analysis.}
For the remainder of this chapter, as well as in
Chapters~\ref{chap:qpbt-combining} and~\ref{chap:qpbt-extraction}, we fix
an admissible parameter tuple $(q,m,d)$
(Definition~\ref{def:admissible}) and a strategy
$\mathcal{S} = (\ket{\psi}, M)$ for the game
$\mathfrak{G}^{\mathrm{Pauli}}_{(q,m,d)}$ that succeeds with probability
at least $1 - \eps$, for some $\eps \geq 0$, where $\ket{\psi}$ is a
bipartite state on registers $\mathsf{A}$ and $\mathsf{B}$ with associated
Hilbert spaces $\hilb_{\mathsf{A}}$ and $\hilb_{\mathsf{B}}$. The symbol
$M$ denotes the measurement operators of both players; which player is
meant is indicated by the register on which the operator acts (for
example, $M^{(\mathrm{Point},W),u}_a \ot \Id_{\mathsf{B}}$ indicates that
$M^{(\mathrm{Point},W),u}_a$ is viewed as an operator acting on register
$\mathsf{A}$).

The analysis requires the measurements of $\mathcal{S}$ to be projective;
the following lemma shows that this may be assumed without loss of
generality.

\begin{lemma}[Reduction to a projective strategy]
	\label{lem:projective-strategy-setup}
	\lean{MIPStarRE.QPBT.exists_projective_padded_strategy}
	\lean{MIPStarRE.QPBT.padWithZeros}
	\lean{MIPStarRE.QPBT.paddedProdShuffle}
	\leanok
	\uses{def:admissible, def:pauli-question-distribution, def:pauli-win-predicate, def:tensor-product-strategy, def:tensor-product-value, def:projective-strategy-general}
	Let $(q,m,d)$ be an admissible parameter tuple
	(Definition~\ref{def:admissible}) and let
	$\mathcal{S} = (\ket{\psi}, M)$ be a strategy for the game
	$\mathfrak{G}^{\mathrm{Pauli}}_{(q,m,d)}$, where $\ket{\psi}$ is a
	bipartite state on registers $\mathsf{A}$ and $\mathsf{B}$. Then there
	exists a projective strategy
	(Definition~\ref{def:projective-strategy-general}) for
	$\mathfrak{G}^{\mathrm{Pauli}}_{(q,m,d)}$ whose success probability
	equals that of $\mathcal{S}$ and whose state is the extension
	\[
		(\ket{\psi} \ot \ket{0\cdots 0} \ot \ket{0\cdots 0})_{\mathsf{A}\mathsf{B}}
	\]
	of $\ket{\psi}$ by finitely many ancilla qubits initialized in the
	state $\ket{0}$, adjoined to each of the two registers.
\end{lemma}
\begin{proof}
	\leanok
	\uses{thm:naimark, lem:naimark-helper, def:measurement-completion,
		lem:bipartite-strategy-block-identities, def:tensor-product-strategy,
		def:tensor-product-value, def:projective-strategy-general}
	Naimark's theorem (Theorem~\ref{thm:naimark}), applied with the
	families of all of the first player's measurements on register
	$\mathsf{A}$ and all of the second player's on register $\mathsf{B}$,
	yields an auxiliary space and a projective sub-measurement replacing
	each family on each register, together with a product auxiliary state,
	such that every correlation
	$\bra{\psi} M^{x}_a \ot M^{y}_b \ket{\psi}$ of $\mathcal{S}$ is
	reproduced on the extended state. Each auxiliary space embeds into a
	register of finitely many qubits; extend each dilated sub-measurement
	to the enlarged register by adding to one fixed outcome both the
	projector onto the orthogonal complement of the embedded subspace and
	the residual projector of the sub-measurement. The outcome operators
	remain mutually orthogonal projectors, now summing to the identity,
	and no correlation changes: the complement projector annihilates the
	extended state, which is supported in the embedded subspace, and so
	does each residual projector, because for each question pair the
	reproduced correlations sum to one --- the original families being
	POVMs (Definition~\ref{def:tensor-product-strategy}) --- which forces
	the projector $\Id$ minus the sum of each dilated sub-measurement to
	have zero mass on the extended state. Finally, since the auxiliary
	state is a product state, a unitary acting on each player's ancilla
	qubits alone carries it to $\ket{0 \cdots 0} \ot \ket{0 \cdots 0}$;
	conjugating the projective measurements by the same unitaries
	preserves projectivity and all correlations. The success probability
	of a strategy is a linear combination of its correlations, weighted by
	the question distribution and the win predicate
	(Definition~\ref{def:tensor-product-value}), and is therefore
	unchanged.
\end{proof}

In view of Lemma~\ref{lem:projective-strategy-setup} we assume from here
on, replacing $\mathcal{S}$ by the projective strategy provided by that
lemma, that all measurements in $\mathcal{S}$ are projective; the
extended state
$(\ket{\psi} \ot \ket{0\cdots 0} \ot \ket{0\cdots 0})_{\mathsf{A}\mathsf{B}}$
is again denoted $\ket{\psi}$. All statements below refer to this
projective strategy on the extended state, and references to ``the
state $\ket{\psi}$ of the strategy under analysis'' refer to this
convention.

The strategy $\mathcal{S}$ contains measurement operators for each of the
question types and question contents of the Pauli basis test
(Definition~\ref{def:pauli-win-predicate}), such as
$\{M^{(\mathrm{Point},W),u}_a\}_{a \in \F_q}$ for all $u \in \F_q^m$.

\begin{definition}[Strategy observables]
	\label{def:strategy-observables}
	\lean{MIPStarRE.QPBT.ProjectiveSetting}
	\lean{MIPStarRE.QPBT.ProjectiveSetting.pointMeas}
	\lean{MIPStarRE.QPBT.ProjectiveSetting.lineMeas}
	\lean{MIPStarRE.QPBT.ProjectiveSetting.pairMeas}
	\lean{MIPStarRE.QPBT.ProjectiveSetting.pairWMeas}
	\lean{MIPStarRE.QPBT.ProjectiveSetting.pauliMeas}
	\lean{MIPStarRE.QPBT.ProjectiveSetting.msMeas}
	\lean{MIPStarRE.QPBT.ProjectiveSetting.wrongFormMass}
	\lean{MIPStarRE.QPBT.ProjectiveSetting.wrongFormMass_alice_le_error}
	\lean{MIPStarRE.QPBT.ProjectiveSetting.wrongFormMass_bob_le_error}
	\lean{MIPStarRE.QPBT.ProjectiveSetting.pointObs}
	\lean{MIPStarRE.QPBT.ProjectiveSetting.pointObs_sq_eq_one}
	\lean{MIPStarRE.QPBT.ProjectiveSetting.pointObs_isHermitian}
	\leanok
	\uses{lem:projective-strategy-setup, def:pauli-win-predicate, def:subfield-trace, def:line, def:line-representative}
	For ``line'' questions we write $\mathrm{Line}$ for either of the
	types $\mathrm{ALine}$ and $\mathrm{DLine}$: a statement about the
	type $\mathrm{Line}$ applies to both types, and the intended type is
	determined by the line description $\ell$. Moreover,
	$\ell$ denotes the description of either an axis-parallel line, given by
	a pair $\ell = (u_0, s)$, or a diagonal line, given by a triple
	$\ell = (u_0, s, v)$ (Definition~\ref{def:line-representative}). For
	the formalization, $M^{(\mathrm{Point},W),u}_a$ denotes the complete
	field-valued measurement obtained by assigning every answer of the wrong
	form to $a=0$. The winning condition bounds the reassigned mass by the
	strategy's rejection probability. For
	every $u \in \F_q^m$, $r \in \F_q$, and $W \in \{X, Z\}$, define
	\begin{equation}\label{eq:qld-strat-obs}
		W^r(u) \,=\, \sum_{a \in \F_q} (-1)^{\tr(ar)} \, M^{(\mathrm{Point},W),u}_a\,.
	\end{equation}
	Since the measurement $\{M^{(\mathrm{Point},W),u}_a\}_{a \in \F_q}$ is
	projective and $(-1)^{\tr(ar)} \in \{\pm 1\}$, the operator $W^r(u)$ is
	an observable with eigenvalues $\pm 1$. It acts on register $\mathsf{A}$
	or $\mathsf{B}$ according to which player's operator
	$M^{(\mathrm{Point},W),u}_a$ denotes.
\end{definition}

\begin{lemma}[Multiplication of point observables]
\label{lem:point-observable-multiplication}
\lean{MIPStarRE.QPBT.ProjectiveSetting.pointObs_mul}
\lean{MIPStarRE.QPBT.ProjectiveSetting.pointObs_zero}
\lean{MIPStarRE.QPBT.ProjectiveSetting.pointMeas_isProjective}
\leanok
\uses{def:strategy-observables}
For a projective strategy, the answer-folded point measurement is projective.
For either player, fixed $W$ and $u$, and all $r,s\in\F_q$, its observables satisfy
\[
W^r(u)W^s(u)=W^{r+s}(u),\qquad W^0(u)=\Id.
\]
These are formalization-only consequences of Equation~\eqref{eq:qld-strat-obs}.
\end{lemma}
\begin{proof}
\leanok
Answer folding preserves projectivity by summing orthogonal projections.
Expanding the product of the two observable sums, orthogonality eliminates
unequal outcomes, and additivity of the trace combines the two phase factors.
At $r=0$, completeness of the point measurement gives the identity.
\end{proof}

\begin{lemma}[Averaging scalar multiples of an operator]
\label{lem:operator-average-scalar-constant}
\lean{MIPStarRE.LDT.averageOperatorOverDistribution_smul_const}
\leanok
Let $D$ be a finite-support distribution on a type $T$, with weights $w(t)$,
let $c:T\to\mathbb C$, and let $A$ be a fixed finite-dimensional operator.
Then
\[
\sum_{t\in\operatorname{supp}(D)} w(t)\bigl(c(t)A\bigr)
=\left(\sum_{t\in\operatorname{supp}(D)}w(t)c(t)\right)A.
\]
This formalization-only identity supplies the scalar averaging rule used in
Fourier inversion for expanded point effects.
\end{lemma}
\begin{proof}
\leanok
Use compatibility of real and complex scalar multiplication and distribute
scalar multiplication over the finite sum.
\end{proof}

The observable $W^r(u)$ is the candidate for the generalized Pauli
observable $\tau^W(r \cdot \mathrm{ind}_m(u))$ that the analysis eventually
extracts from the strategy.

\begin{lemma}[Winning implications]
	\label{lem:qld-win-implications}
	\lean{MIPStarRE.QPBT.win_cons, MIPStarRE.QPBT.win_cons_approx}
	\lean{MIPStarRE.QPBT.win_low_degree, MIPStarRE.QPBT.win_low_degree_approx}
	\lean{MIPStarRE.QPBT.win_pauli_basis_cons, MIPStarRE.QPBT.win_pauli_basis_cons_approx}
	\lean{MIPStarRE.QPBT.win_comm, MIPStarRE.QPBT.win_comm_approx}
	\lean{MIPStarRE.QPBT.win_comm_cons, MIPStarRE.QPBT.win_comm_cons_approx}
	\lean{MIPStarRE.QPBT.win_magic_square}
	\lean{MIPStarRE.QPBT.win_ms_cons, MIPStarRE.QPBT.win_ms_cons_approx}
	% The conditional sampling of items 4-7: the uniform laws on commuting
	% and on anticommuting tuples, together with their total mass one.
	\lean{MIPStarRE.QPBT.anticommTupleDist, MIPStarRE.QPBT.commTupleDist}
	\lean{MIPStarRE.QPBT.anticommTupleDist_isProbability}
	\lean{MIPStarRE.QPBT.commTupleDist_isProbability}
	\leanok
	\uses{lem:projective-strategy-setup, def:anticommuting-tuple, def:pauli-question-distribution, def:pauli-win-predicate, def:line-point-dist, def:consistency, def:povm-distance, def:ms-game, def:tensor-product-value, def:low-degree-encoding, def:ideg-deg-polynomials}
	Assume that $\mathcal{S}$ succeeds in $\mathfrak{G}^{\mathrm{Pauli}}_{(q,m,d)}$
	with probability at least $1 - \eps$, for some $\eps \geq 0$. Then the
	following hold, with respect to the state $\ket{\psi}$ of the strategy
	under analysis (Lemma~\ref{lem:projective-strategy-setup} and the
	convention following it).
	\begin{enumerate}
		\item \label{enu:qld-cons} (\textbf{Consistency check}) On average
		over $(\mathsf{t},x)$ sampled from the question distribution of
		$\mathfrak{G}^{\mathrm{Pauli}}_{(q,m,d)}$
		(Definition~\ref{def:pauli-question-distribution}), we have
		$M^{\mathsf{t},x}_a \ot \Id_{\mathsf{B}} \simeq_\eps \Id_{\mathsf{A}} \ot M^{\mathsf{t},x}_a$.
		\item \label{enu:qld-ld} (\textbf{Low-degree check}) For all
		$W \in \{X,Z\}$, on average over $(\ell,u)$ drawn from the
		line-point distribution (Definition~\ref{def:line-point-dist}), we
		have
		% The source states this item with answer summation over F_q only,
		% implicitly reading the evaluation classes of the line measurement
		% as a complete measurement; see the paragraph after the item list
		% and the remark rem:qld-win-implications-typos.
		$M^{(\mathrm{Line},W),\ell}_{[\mathrm{eval}_u(\cdot)=a]} \ot \Id_{\mathsf{B}} \simeq_\eps \Id_{\mathsf{A}} \ot M^{(\mathrm{Point},W),u}_a$,
		where the answer summation is over the completed outcome set
		$a \in \F_q \cup \{\bot\}$ described after the item list.
		\item \label{enu:qld-pauli-basis-cons} (\textbf{Pauli basis
		consistency check}) For all $W \in \{X,Z\}$, on average over
		uniformly random $u \in \F_q^m$, we have
		$M^{(\mathrm{Point},W),u}_a \ot \Id_{\mathsf{B}} \simeq_\eps \Id_{\mathsf{A}} \ot M^{(\mathrm{Pauli},W)}_{[g_h(u)=a]}$,
		where
		$M^{(\mathrm{Pauli},W)}_{[g_h(u)=a]} = \sum_{h \in \F_q^{M} \colon g_h(u)=a} M^{(\mathrm{Pauli},W)}_h$
		and $g_h$ is the low-degree encoding of $h$
		(see Definition~\ref{def:low-degree-encoding}).
		\item \label{enu:qld-comm} (\textbf{Commutation check}) For all
		$W \in \{X,Z\}$, on average over commuting
		$\omega = (u_X,u_Z,r_X,r_Z)$ (that is, over a uniformly random
		$\omega$ conditioned on being commuting), we have
		$M^{(\mathrm{Pair},W),\omega}_a \ot \Id_{\mathsf{B}} \simeq_\eps \Id_{\mathsf{A}} \ot M^{\mathrm{Pair},\omega}_{[\beta_W=a]}$,
		where the answer summation is over $a \in \F_2$.
		\item \label{enu:qld-comm-cons} (\textbf{Commutation consistency
		check}) For all $W \in \{X,Z\}$, on average over commuting
		$\omega$, we have
		$M^{(\mathrm{Point},W),u_W}_{[\tr(\cdot\,r_W)=a]} \ot \Id_{\mathsf{B}} \simeq_\eps \Id_{\mathsf{A}} \ot M^{(\mathrm{Pair},W),\omega}_a$,
		where the answer summation is over $a \in \F_2$.
		\item \label{enu:qld-ms} (\textbf{Magic square check}) For any
		anticommuting $\omega$, let $\Lambda_\omega$ be the value of the
		strategy
		$\big( \ket{\psi}, \{M^{\mathrm{Constraint}_i,\omega}_{\alpha_1,\alpha_2,\alpha_3}\} \cup \{M^{\mathrm{Variable}_j,\omega}_a\} \big)$
		in the Magic Square game $\mathfrak{G}^{\mathrm{MS}}$
		(Definition~\ref{def:ms-game}). Then
		$\E_\omega \Lambda_\omega = 1 - O(\eps)$, where the expectation is
		over a uniformly random $\omega$ conditioned on being
		anticommuting.
		\item \label{enu:qld-ms-cons} (\textbf{Magic square consistency
		check}) For all $W \in \{X,Z\}$, on average over anticommuting
		$\omega$,
		\begin{gather*}
			M^{(\mathrm{Point},X),u_X}_{[\tr(\cdot\,r_X)=a]} \ot \Id_{\mathsf{B}} \,\simeq_\eps\, \Id_{\mathsf{A}} \ot M^{\mathrm{Variable}_1,\omega}_a\,,\\
			M^{(\mathrm{Point},Z),u_Z}_{[\tr(\cdot\,r_Z)=a]} \ot \Id_{\mathsf{B}} \,\simeq_\eps\, \Id_{\mathsf{A}} \ot M^{\mathrm{Variable}_5,\omega}_a\,,
		\end{gather*}
		where the answer summations are over $\F_2$.
	\end{enumerate}
	Here
	$M^{(\mathrm{Point},W),u}_{[\tr(\cdot\,r)=a]} = \sum_{a' \in \F_q \colon \tr(a'r)=a} M^{(\mathrm{Point},W),u}_{a'}$
	for $a \in \F_2$; in item~\ref{enu:qld-ld} the two families are the
	following POVMs indexed by $\F_q \cup \{\bot\}$, as the relation
	$\simeq_\eps$ of Definition~\ref{def:consistency} requires: on the left,
	the completed family of evaluation classes of
	Definition~\ref{def:ideg-deg-polynomials}, whose class
	$M^{(\mathrm{Line},W),\ell}_{[\mathrm{eval}_u(\cdot)=\bot]}$ carries the
	line answers with no evaluation at $u$ and is nonzero only for a
	zero-direction line, where the classes indexed by $a \in \F_q$ form
	only a sub-measurement; on the right, the point measurement completed
	by the zero operator $M^{(\mathrm{Point},W),u}_{\bot} = 0$; the answer
	to a question of type $\mathrm{Pair}$ is a
	pair $(\beta_X,\beta_Z) \in \F_2 \times \F_2$, and
	$M^{\mathrm{Pair},\omega}_{[\beta_W=a]}$ denotes its coarse-graining
	onto the $W$-component. Furthermore, all of the approximations above
	also hold with $\approx_\eps$ in place of $\simeq_\eps$, and with the
	tensor factors interchanged.  The corrected conditioning event,
	zero-direction completion, and unrestricted parameter proof are documented
	in \cite{gap:qpbt_win-implications-corrections}.
\end{lemma}
\begin{proof}
	\leanok
	\uses{fact:omega-anticomm-prob, fact:agreement, def:admissible, def:pauli-question-distribution, def:pauli-win-predicate, def:ld-win-predicate, def:ideg-deg-polynomials, def:consistency, lem:qld-reversed-consistency, lem:qld-distance-interchange, lem:qld-placed-consistency-distance}
	The question types are sampled according to the type graph of the Pauli
	basis test (Definition~\ref{def:pauli-question-distribution}), which has
	constant size, so each subtest of the win predicate
	(Definition~\ref{def:pauli-win-predicate}) is executed with probability
	bounded below by a universal constant.
	Items~\ref{enu:qld-cons},~\ref{enu:qld-ld}
	and~\ref{enu:qld-pauli-basis-cons} then follow directly from the
	corresponding subtests and the assumption that $\mathcal{S}$ succeeds
	with probability at least $1-\eps$, where item~\ref{enu:qld-ld} follows
	from the consistency subtest of the low individual degree game executed
	inside the Pauli basis test. In more detail for item~\ref{enu:qld-ld}:
	the disagreement mass of the two completed families is rejection mass
	of the line-versus-point clauses
	(Definition~\ref{def:ld-win-predicate}, invoked by the low-degree
	check of Definition~\ref{def:pauli-win-predicate}). An outcome pair
	$(a,b)$ with $a \neq b$ in $\F_q$ is a line answer evaluating at $u$
	to $a$ against the point answer $b$, which the clause rejects; a pair
	$(\bot,b)$ is a line answer with no evaluation at $u$, which the
	universally quantified clause rejects against every point answer $b$
	(Remark~\ref{rem:ld-win-zero-direction} and
	Remark~\ref{rem:deg-line-representatives}); and the pairs $(a,\bot)$
	contribute no mass, since the class
	$M^{(\mathrm{Point},W),u}_{\bot}$ is the zero operator. The subtests underlying
	items~\ref{enu:qld-comm} and~\ref{enu:qld-comm-cons} accept vacuously
	unless the sampled tuple $\omega$ is commuting, and those underlying
	items~\ref{enu:qld-ms} and~\ref{enu:qld-ms-cons} accept vacuously
	unless $\omega$ is anticommuting. Since the parameter tuple $(q,m,d)$
	is admissible (Definition~\ref{def:admissible}), $m$ divides $q$, so
	$1 \leq m \leq q$, and Lemma~\ref{fact:omega-anticomm-prob} shows that
	a uniformly sampled tuple $\omega$ is commuting with probability at
	least $\frac12$ and anticommuting with probability at least $2^{-4}$.
	Conditioning on either event therefore increases the failure
	probability of the corresponding subtest by a factor of at most $2^4$,
	and items~\ref{enu:qld-comm},~\ref{enu:qld-comm-cons}
	and~\ref{enu:qld-ms-cons} follow from the corresponding subtests.
	Item~\ref{enu:qld-ms} follows in the same way, because the Magic Square
	check is executed with constant probability over the choice of a pair
	of questions and the anticommuting event has probability at least
	$2^{-4}$.  This replacement for the source's circular parameter reduction
	is documented in \cite{gap:qpbt_win-implications-corrections}. Finally, that all approximations
	also hold with $\approx_\eps$ in place of $\simeq_\eps$ is immediate
	from the quantitative placed form of Lemma~\ref{fact:agreement}
	in Lemma~\ref{lem:qld-placed-consistency-distance}.
	The forms with the tensor factors interchanged are the same subtests
	read on the reversed ordered question pair. That pair is again sampled
	by the type graph, whose edges are unordered
	(Definition~\ref{def:pauli-question-distribution}), and the win
	predicate (Definition~\ref{def:pauli-win-predicate}) lists every clause
	together with its mirror image, so exchanging the two questions and the
	two answers leaves it unchanged. Interchanging the two tensor factors
	relabels the joint coordinates: it exchanges the two placements of a pair
	of operators and carries $\ket{\psi}$ to the interchanged state, so the
	state-dependent distance of the two families is unchanged.
	The reversed consistency estimates are recorded in
	Lemma~\ref{lem:qld-reversed-consistency}; the distance identity used
	after applying Lemma~\ref{fact:agreement} is
	Lemma~\ref{lem:qld-distance-interchange}.
\end{proof}

\begin{lemma}[Placed consistency-to-distance estimates]
	\label{lem:qld-placed-consistency-distance}
	\lean{MIPStarRE.QPBT.WinImplications.opFamilyDistSq_placed_le_two_mul_consistencyDefect}
	\lean{MIPStarRE.QPBT.WinImplications.opFamilyDistSq_placed_le_of_consistencyDefect_le}
	\leanok
	\uses{def:consistency, def:povm-distance}
	Let $X$ and $\mathcal A$ be finite question and outcome sets, and let
	$\mu$ be a nonnegative weight on $X$. For each $x\in X$, let
	$\{A_a^x\}_{a\in\mathcal A}$ and $\{B_a^x\}_{a\in\mathcal A}$ be POVMs
	on finite-dimensional complex coordinate spaces $\mathcal H_A$ and
	$\mathcal H_B$, respectively. For any vector
	$\psi\in\mathcal H_A\otimes\mathcal H_B$, set
	\[
	D_\mu=\sum_x\mu(x)\sum_{a\in\mathcal A}
	 \big\|(A_a^x\otimes\Id_B-\Id_A\otimes B_a^x)\psi\big\|^2,
	\qquad
	\mathrm{Cons}_\mu=\sum_x\mu(x)\sum_{a\ne b}
	 \langle\psi|A_a^x\otimes B_b^x|\psi\rangle.
	\]
	Then $D_\mu\leq 2\,\mathrm{Cons}_\mu$. Consequently, for every
	$c\in\mathbb R$, if $\mathrm{Cons}_\mu\leq c$, then $D_\mu\leq 2c$.
	Neither $\mu$ nor $\psi$ needs to be normalized, and the POVMs need
	not be projective. These are formalization-only quantitative auxiliaries
	for the final sentence of Lemma~\ref{lem:qld-win-implications}.
\end{lemma}
\begin{proof}
	\leanok
	\uses{fact:agreement}
	Apply the factor-two consistency-to-distance inequality of
	Lemma~\ref{fact:agreement} to the POVMs
	$\{A_a^x\otimes\Id_B\}$ and $\{\Id_A\otimes B_a^x\}$ on the joint
	space. Their off-diagonal products are $A_a^x\otimes B_b^x$,
	giving the first inequality. The second follows from
	$D_\mu\leq 2\,\mathrm{Cons}_\mu\leq 2c$.
\end{proof}

\begin{lemma}[Distance under tensor interchange]
	\label{lem:qld-distance-interchange}
	\lean{MIPStarRE.QPBT.WinImplications.opFamilyDistSq_swappedState}
	\leanok
	\uses{def:povm-distance}
	Let $\mathcal H_A,\mathcal H_B$ be finite-dimensional complex coordinate
	spaces, $\mathcal A$ a finite outcome set, and $\mu$ a finitely supported
	nonnegative weight on a question set $X$. Let $M_a^x$ and $N_a^x$ be
	arbitrary operators on $\mathcal H_A$ and $\mathcal H_B$, respectively.
	For any vector $\psi\in\mathcal H_A\otimes\mathcal H_B$, let
	$U:\mathcal H_A\otimes\mathcal H_B\to\mathcal H_B\otimes\mathcal H_A$
	be the coordinate interchange $U(v\otimes w)=w\otimes v$. Then
	\[
	\sum_x\mu(x)\sum_{a\in\mathcal A}
	 \big\|(N_a^x\otimes\Id_A-\Id_B\otimes M_a^x)U\psi\big\|^2
	=
	\sum_x\mu(x)\sum_{a\in\mathcal A}
	 \big\|(M_a^x\otimes\Id_B-\Id_A\otimes N_a^x)\psi\big\|^2.
	\]
	This is a formalization-only auxiliary identity for the final sentence
	of Lemma~\ref{lem:qld-win-implications}; neither normalization nor
	positivity of the operators is required.
\end{lemma}
\begin{proof}
	\leanok
	Conjugating by $U$ exchanges the tensor placements. The operator
	difference on the left becomes the negative of that on the right.
	Coordinate interchange and negation preserve the norm, so the identity
	holds term by term in both sums.
\end{proof}

\begin{lemma}[Reversed-orientation consistency estimates]
	\label{lem:qld-reversed-consistency}
	\lean{MIPStarRE.QPBT.WinImplications.win_low_degree_interchanged_proof}
	\lean{MIPStarRE.QPBT.WinImplications.win_pauli_basis_cons_interchanged_proof}
	\lean{MIPStarRE.QPBT.WinImplications.win_comm_interchanged_proof}
	\lean{MIPStarRE.QPBT.WinImplications.win_comm_cons_interchanged_proof}
	\lean{MIPStarRE.QPBT.WinImplications.win_ms_cons_interchanged_proof}
	\leanok
	\uses{lem:projective-strategy-setup, def:admissible, def:consistency,
		def:line-point-dist, def:low-degree-encoding, def:ideg-deg-polynomials}
	For each of Items~\ref{enu:qld-ld}, \ref{enu:qld-pauli-basis-cons},
	\ref{enu:qld-comm}, \ref{enu:qld-comm-cons}, and \ref{enu:qld-ms-cons},
	there is a universal real constant $C\geq1$ with the following property.
	For every admissible parameter tuple and projective strategy succeeding
	in the Pauli basis test with probability at least $1-\eps$, where
	$\eps\geq0$, the reversed consistency defect is at most $C\eps$,
	for each $W\in\{X,Z\}$.
	More precisely, write $L_a^x$ and $R_a^x$ for the local measurement
	families on the left and right of the corresponding item, and $\mu$
	for its sampling distribution. With superscripts $A,B$ indicating
	which player's measurement is used, the assertion is
	\[
	\E_{x\sim\mu}\sum_{a\ne b}
	 \langle\psi|R_a^{x,A}\otimes L_b^{x,B}|\psi\rangle
	\leq C\eps.
	\]
	The low-degree item uses the completed outcome set
	$\F_q\cup\{\bot\}$ and line-point distribution described above;
	the Pauli basis item uses uniform points and outcomes in $\F_q$;
	the remaining three items use binary outcomes and the respective
	commuting or anticommuting conditional tuple distribution.
	These auxiliary estimates give the reversed forms of precisely those
	five items, without assuming that the strategy is symmetric.
\end{lemma}
\begin{proof}
	\leanok
	\uses{def:pauli-question-distribution, def:pauli-win-predicate,
		def:ld-win-predicate, fact:omega-anticomm-prob}
	Each disagreement event is rejected by the clause for the reversed
	ordered question pair. Averaging over that clause bounds its mass by
	the total rejection probability times the constant number of type-graph
	edges. For the last three items, conditioning costs at most $2$, $2$,
	and $16$, respectively, by Lemma~\ref{fact:omega-anticomm-prob}.
	For the low-degree item the completed outcome $\bot$ is treated by
	the same rejection clause, with zero mass on the point side.
\end{proof}

\begin{lemma}[Consistency and value part of Lemma~\ref{lem:qld-win-implications}]
	\label{lem:qld-win-implications-exact}
	\lean{MIPStarRE.QPBT.win_cons, MIPStarRE.QPBT.win_low_degree}
	\lean{MIPStarRE.QPBT.win_pauli_basis_cons, MIPStarRE.QPBT.win_comm}
	\lean{MIPStarRE.QPBT.win_comm_cons, MIPStarRE.QPBT.win_magic_square}
	\lean{MIPStarRE.QPBT.win_ms_cons}
	\leanok
	\uses{lem:qld-win-implications}
	Under the hypotheses and notation of
	Lemma~\ref{lem:qld-win-implications}, each of its seven numbered conclusions
	holds in the displayed tensor order: Items~\ref{enu:qld-cons}--\ref{enu:qld-comm-cons}
	and~\ref{enu:qld-ms-cons} hold with the consistency relation
	$\simeq_\eps$, and Item~\ref{enu:qld-ms} gives the stated average Magic
	Square value. This assertion does not include the operator-distance or
	tensor-interchanged conclusions in the final sentence of that lemma.
\end{lemma}
\begin{proof}
	\leanok
	\uses{fact:omega-anticomm-prob, lem:graph-distribution-mass,
		def:admissible,
		def:pauli-question-distribution, def:pauli-win-predicate,
		def:ld-win-predicate, def:line-point-dist,
		def:ideg-deg-polynomials, def:consistency, def:ms-game,
		def:tensor-product-value}
	For each consistency assertion, rewrite the defect as the mass of unequal
	coarse-grained labels. Every such mismatch violates the corresponding clause
	of the Pauli basis test, so its mass is bounded by rejection on the associated
	edge of the finite type graph; for the low-degree check, first split the
	line-point average into its axis-parallel and diagonal parts. Uniform edge
	sampling bounds each resulting rejection average by a universal constant
	times $\eps$. For the commuting and anticommuting items, conditioning costs
	factors at most $2$ and $2^4$, respectively, by
	Lemma~\ref{fact:omega-anticomm-prob}. Finally, identify one minus the value of
	the induced Magic Square strategy with its average rejection mass, apply the
	same anticommuting and fixed-edge bounds, and average over the Magic Square
	question pairs to obtain Item~\ref{enu:qld-ms}.
\end{proof}

\begin{lemma}[Magic Square anticommutation on the original state]
	\label{lem:qld-ms-anticommutator-original-state}
	\lean{MIPStarRE.QPBT.WinImplications.msVarObs_anticommutator_le}
	\lean{MIPStarRE.QPBT.WinImplications.msVarObsA_anticommutator_le}
	\leanok
	\uses{def:ms-game, def:tensor-product-value, def:ms-variable-agreement}
	This auxiliary estimate isolates the one-player anticommutation used in
	the proof of Lemma~\ref{lem:qld-win-implications-obs}.
	Let $\mathcal S=(\ket{\psi},A,B)$ be a Magic Square strategy of value
	at least $1-\eta$, with $\eta\geq0$. For $j\in\{1,5\}$, let
	$V_j^B=\hat B^{\mathsf{Variable}_j}_0-\hat B^{\mathsf{Variable}_j}_1$,
	where the binary effects are those of
	Definition~\ref{def:ms-variable-agreement}. Then
	\[
	\big\|(\Id_A\otimes(V_1^B V_5^B+V_5^B V_1^B))\ket{\psi}\big\|^2
	\leq 1183680\,\eta.
	\]
	The same bound holds for Alice's binary variable observables, placed
	on the first tensor factor.
\end{lemma}
\begin{proof}
	\leanok
	\uses{thm:naimark, def:ms-game}
	On the questionwise projective dilation, the constraint--variable
	correlations and the Magic Square parity relations bound the norm of
	each player's anticommutator by $624\sqrt\eta$. Compression to the
	original space bounds the squared anticommutator norm by three times
	this squared norm plus $15552\eta$.
	The dilation embedding preserves norms, and an inflated operator acts
	on the embedded state as the original operator acts on $\ket{\psi}$.
	Thus the original-state bound is
	$(3\cdot624^2+15552)\eta=1183680\eta$.
\end{proof}

\begin{lemma}[Point commutation averaged over commuting tuples]
	\label{lem:qld-point-obs-commuting-estimate}
	\lean{MIPStarRE.QPBT.WinImplications.exists_pointObs_commutator_comm_le_alice}
	\leanok
	\uses{lem:projective-strategy-setup, def:strategy-observables,
		def:anticommuting-tuple}
	This auxiliary estimate is the commuting part of the proof of
	Lemma~\ref{lem:qld-win-implications-obs}.
	There is a universal real constant $C\geq1$ such that, for every
	admissible parameter tuple and projective strategy succeeding in the
	Pauli basis test with probability at least $1-\eps$, where $\eps\geq0$,
	\[
	\E_{\omega\text{ commuting}}
	\big\|\big((X^{r_X}(u_X)Z^{r_Z}(u_Z)
	 -Z^{r_Z}(u_Z)X^{r_X}(u_X))\otimes\Id_B\big)\ket{\psi}\big\|^2
	\leq C(\eps+\sqrt\eps).
	\]
	The expectation is uniform conditional on $\gamma(\omega)=0$.
\end{lemma}
\begin{proof}
	\leanok
	\uses{lem:qld-win-implications-exact, fact:triangle-for-simeq,
		fact:agreement, lem:commutation-analysis, def:strategy-observables}
	Use the point--$\mathrm{Pair}/W$ consistency, the
	$\mathrm{Pair}/W$ self-consistency, and the
	$\mathrm{Pair}/W$--$\mathrm{Pair}$ marginal consistency.
	If their defects are $c_1,c_2,c_3$, respectively, then
	Lemmas~\ref{fact:triangle-for-simeq} and~\ref{fact:agreement} bound
	the squared point--$\mathrm{Pair}$ marginal distance by
	$2(c_1+2\sqrt{c_2+c_3})$. Each $c_i$ is $O(\eps)$ by
	Lemma~\ref{lem:qld-win-implications-exact}, so this is
	$O(\eps+\sqrt\eps)$. The $\mathrm{Pair}$ measurement is projective;
	Lemma~\ref{lem:commutation-analysis} therefore gives the same order
	for the projection commutator. Expanding
	$W^r(u)=\Id-2M^{(\mathrm{Point},W),u}_{[\tr(\cdot\,r)=1]}$
	multiplies that commutator by $4$, preserving the order of the bound.
\end{proof}

\begin{lemma}[Winning implications for observables]
	\label{lem:qld-win-implications-obs}
	\lean{MIPStarRE.QPBT.pointObs_self_consistent}
	\lean{MIPStarRE.QPBT.pointObs_twisted_commutation}
	\lean{MIPStarRE.QPBT.pointObs_twisted_commutation_interchanged}
	\leanok
	% The source labels this lemma `len:qld-win-implications-obs'; the prefix
	% `len' is a typo for `lem'. The label above is the corrected form.
	\uses{def:strategy-observables, def:approx-question-indexed-operators, def:anticommuting-tuple, lem:projective-strategy-setup}
	The following hold for the observables $W^r(u)$ of
	Definition~\ref{def:strategy-observables}, with respect to the state
	$\ket{\psi}$ of the strategy under analysis
	(Lemma~\ref{lem:projective-strategy-setup} and the convention
	following it).
	\begin{enumerate}
		\item For all $W \in \{X,Z\}$ and $r \in \F_q$, on average over
		uniformly random $u \in \F_q^m$,
		\begin{equation}\label{eq:pts-obs-consistency}
			W^r(u) \ot \Id_{\mathsf{B}} \,\approx_\eps\, \Id_{\mathsf{A}} \ot W^r(u)\,.
		\end{equation}
		\item On average over uniformly random
		$\omega = (u_X,u_Z,r_X,r_Z) \in (\F_q^m)^2 \times (\F_q)^2$,
		\begin{equation}\label{eq:pts-obs-commutation}
			X^{r_X}(u_X)\, Z^{r_Z}(u_Z) \ot \Id_{\mathsf{B}}
			\,\approx_{\sqrt{\eps}}\,
			(-1)^{\gamma(\omega)}\, Z^{r_Z}(u_Z)\, X^{r_X}(u_X) \ot \Id_{\mathsf{B}}\,,
		\end{equation}
		where
		$\gamma(\omega) = \tr((r_X \cdot \mathrm{ind}_m(u_X)) \cdot (r_Z \cdot \mathrm{ind}_m(u_Z)))$
		as in Definition~\ref{def:anticommuting-tuple}.
	\end{enumerate}
	Furthermore, all of the approximations above also hold with the tensor
	factors interchanged.
\end{lemma}
\begin{proof}
	\leanok
	\uses{lem:qld-win-implications-exact, lem:povm-to-obs,
		fact:data-processing, fact:agreement, fact:add-a-proj,
		lem:qld-ms-anticommutator-original-state,
		lem:qld-point-obs-commuting-estimate, lem:qld-reversed-consistency,
		lem:qld-distance-interchange, def:strategy-observables}
	For~\eqref{eq:pts-obs-consistency}: item~\ref{enu:qld-cons} of
	Lemma~\ref{lem:qld-win-implications} and the data processing inequality
	(Lemma~\ref{fact:data-processing}), using that the type
	$(\mathrm{Point},W)$ is sampled with constant probability, give the
	self-consistency of the coarse-grained measurements
	$\{M^{(\mathrm{Point},W),u}_{[\tr(\cdot\,r)=a]}\}_{a \in \F_2}$ on
	average over $u$; Lemma~\ref{lem:povm-to-obs}, applied with
	$\mathcal{A} = \F_2$ and $\alpha_a = (-1)^a$, converts this
	to~\eqref{eq:pts-obs-consistency}. For the commuting case of
	\eqref{eq:pts-obs-commutation},
	Lemma~\ref{lem:qld-point-obs-commuting-estimate} gives a conditional
	average bounded by $C_c(\eps+\sqrt\eps)$.
	For the anticommuting case, item~\ref{enu:qld-ms-cons} of
	Lemma~\ref{lem:qld-win-implications} together with
	Lemma~\ref{lem:povm-to-obs} relates $X^{r_X}(u_X)$ and $Z^{r_Z}(u_Z)$
	to the Magic Square observables
	$M^{\mathrm{Variable}_j,\omega} = M^{\mathrm{Variable}_j,\omega}_0 - M^{\mathrm{Variable}_j,\omega}_1$
	for $j \in \{1,5\}$. Apply
	Lemma~\ref{lem:qld-ms-anticommutator-original-state} to the induced
	Magic Square strategy at each anticommuting $\omega$, with
	$\eta=\eps_\omega=1-\Lambda_\omega$. This gives a squared
	anticommutator norm at most $1183680\eps_\omega$ on $\ket{\psi}$.
	Item~\ref{enu:qld-ms} bounds $\E\eps_\omega$ by $O(\eps)$, so
	direct averaging gives an $O(\eps)$ bound for Bob's variable
	anticommutator. Transferring both products across the tensor factors
	using the point--variable consistency estimates and the reflection
	identities gives a conditional point-anticommutator bound $C_a\eps$.
	The uniform average is bounded by the sum of the two conditional
	averages. For $0\leq\eps\leq1$, use $\eps\leq\sqrt\eps$ to
	obtain $(2C_c+C_a)\sqrt\eps$. For $\eps\geq1$, each product of
	reflections sends $\ket{\psi}$ to a unit vector, so the squared
	distance is at most $4\leq4\sqrt\eps$.
	Finally, interchange the tensor factors and use the reversed
	consistency estimates of Lemma~\ref{lem:qld-reversed-consistency}
	and Alice's original-state Magic Square bound. Coordinate interchange
	preserves all the state-dependent norms, giving the final assertion.
\end{proof}

\begin{remark}[Source corrections]
	\label{rem:qld-win-implications-typos}
	Two typographical slips in the source are corrected above. In
	item~\ref{enu:qld-ms} of Lemma~\ref{lem:qld-win-implications} the
	source writes the expectation over $\omega$ subject to the condition
	$\tr((r_X \cdot u_X) \cdot (r_Z \cdot u_Z)) \neq 0$, omitting the map
	$\mathrm{ind}_m$; the intended condition, consistent with the
	surrounding text and with Definition~\ref{def:anticommuting-tuple}, is
	$\gamma(\omega) \neq 0$, that is, $\omega$ anticommuting, and this is
	the form stated above. In the source proof of
	Lemma~\ref{lem:qld-win-implications-obs}, one display reads
	$Z^{r_X}(u_X)\,X^{r_Z}(u_Z)$ where $Z^{r_Z}(u_Z)\,X^{r_X}(u_X)$ is
	meant.
	% A third slip in the source, the label prefix `len' in place of `lem'
	% on the lemma corresponding to lem:qld-win-implications-obs, is
	% recorded in the comment at that lemma; no alias label is defined.
	In addition, the source proof of Lemma~\ref{lem:qld-win-implications}
	restricts attention to the case $6md \leq q$ and disposes of the
	remaining parameters by an appeal to Theorem~\ref{thm:pauli}, the
	theorem whose proof the lemma is part of; that argument does not
	establish the unrestricted statement of the lemma, and for $d = 0$ the
	probability bound it invokes is false
	(Remark~\ref{rem:omega-anticomm-prob-correction}). The proof given
	above instead uses the lower bounds of
	Lemma~\ref{fact:omega-anticomm-prob} on the probabilities of the
	commuting and anticommuting events, which are universal constants for
	every admissible parameter tuple, so no restriction on the parameters
	is needed.
	Finally, the source states item~\ref{enu:qld-ld} of
	Lemma~\ref{lem:qld-win-implications} with answer summation over
	$a \in \F_q$ only, implicitly reading the evaluation classes of the
	line measurement as a complete measurement. Under the coefficient
	reading of Definition~\ref{def:ideg-deg-polynomials}
	(Remark~\ref{rem:deg-line-representatives}) the classes indexed by
	$a \in \F_q$ form only a sub-measurement when the line has zero
	direction --- a case the question distribution produces with positive
	probability (Remark~\ref{rem:ld-win-zero-direction}) --- while the
	relation $\simeq_\eps$ of Definition~\ref{def:consistency} is defined
	for POVMs. Item~\ref{enu:qld-ld} therefore completes both families by
	an explicit $\bot$ outcome; the same implicit completeness assumption
	occurs in the source proof of the statement corresponding to
	Lemma~\ref{lem:qld-comm-line-cons}, whose proof below invokes the
	completed families instead.  These corrections are collected in
	\cite{gap:qpbt_win-implications-corrections}.
\end{remark}

\section{Expanding the Hilbert space and defining commuting observables}
\label{sec:expanding}

The next step of the analysis adjoins ancillary EPR registers to the
players' state and lifts the strategy observables and measurements to the
expanded state. The point of the construction is that the ancilla factor of
each lifted observable carries the \emph{exact} Pauli commutation sign, so
that the signs cancel and the lifted observables approximately commute for
every tuple $\omega$.

\begin{definition}[Expanded state]
	\label{def:expanded-state}
	\lean{MIPStarRE.QPBT.SixReg, MIPStarRE.QPBT.ProjectiveSetting.psiHat}
	\leanok
	\uses{lem:projective-strategy-setup, def:EPR}
	Recall that $M = 2^m$. Introduce registers $\mathsf{A}', \mathsf{A}'',
	\mathsf{B}', \mathsf{B}''$, each with Hilbert space isomorphic to
	$(\C^q)^{\ot M}$, and define the state
	\begin{equation}\label{eq:def-psihat}
		\ket{\hat{\psi}}_{\mathsf{A}\mathsf{A}'\mathsf{A}''\mathsf{B}\mathsf{B}'\mathsf{B}''}
		\,=\, (\ket{\psi})_{\mathsf{A}\mathsf{B}}
		\ot \ket{\mathrm{EPR}_q}^{\ot M}_{\mathsf{A}'\mathsf{A}''}
		\ot \ket{\mathrm{EPR}_q}^{\ot M}_{\mathsf{B}'\mathsf{B}''}\,,
	\end{equation}
	where $\ket{\psi}$ is the extended strategy state of the strategy
	under analysis (Lemma~\ref{lem:projective-strategy-setup} and the
	convention following it) and
	$\ket{\mathrm{EPR}_q}$ is the maximally entangled state on
	$\C^q \ot \C^q$ (Definition~\ref{def:EPR}). For the remainder of this
	chapter, as well as in Chapters~\ref{chap:qpbt-combining}
	and~\ref{chap:qpbt-extraction}, all approximations $\approx$ and
	$\simeq$ are evaluated with respect to the state $\ket{\hat{\psi}}$,
	unless stated otherwise.
\end{definition}

\begin{definition}[Expanded observables]
	\label{def:expanded-observables}
	\lean{MIPStarRE.QPBT.ProjectiveSetting.expObs}
	\leanok
	\uses{def:strategy-observables, def:expanded-state, lem:twisted-commutation}
	For all $W \in \{X,Z\}$, $r \in \F_q$, and $u \in \F_q^m$, define the
	observable
	\begin{equation}\label{eq:lc-23}
		\hat{W}^r(u) \,=\, W^r(u) \ot \tau^W(r \cdot \mathrm{ind}_m(u))\,,
	\end{equation}
	where $W^r(u)$ is the strategy observable of~\eqref{eq:qld-strat-obs}
	and $\tau^W(r \cdot \mathrm{ind}_m(u))$ is the generalized Pauli
	observable acting on $(\C^q)^{\ot M}$, introduced in
	Chapter~\ref{chap:qpbt-algebra}
	(see Lemma~\ref{lem:twisted-commutation}). As a product
	of two observables with eigenvalues $\pm 1$ acting on disjoint
	registers, $\hat{W}^r(u)$ is itself an observable with eigenvalues
	$\pm 1$. If $W^r(u)$ is viewed as acting on register $\mathsf{A}$, then
	$\hat{W}^r(u)$ is viewed as acting on registers
	$\mathsf{A}\mathsf{A}'$, with the $\tau^W$ factor acting on
	$\mathsf{A}'$. The registers are generally indicated by subscripts, as
	in $(\hat{W}^r(u))_{\mathsf{A}\mathsf{A}'}$ or
	$(\hat{W}^r(u))_{\mathsf{B}\mathsf{A}''}$, the latter having the
	strategy factor on $\mathsf{B}$ and the Pauli factor on
	$\mathsf{A}''$.
\end{definition}

\begin{definition}[Expanded point measurements]
	\label{def:expanded-point-measurement}
	\lean{MIPStarRE.QPBT.ProjectiveSetting.expPointOp}
	\lean{MIPStarRE.QPBT.ProjectiveSetting.tauPointProj}
	\lean{MIPStarRE.QPBT.ProjectiveSetting.pointMeasExp}
	\leanok
	\uses{def:expanded-observables, lem:pauli-observable-expansion, def:low-degree-encoding, def:subfield-trace}
	For $W \in \{X,Z\}$, $u \in \F_q^m$ and $a \in \F_q$, define
	\[
		\hat{M}^{(\mathrm{Point},W),u}_a \,=\, \E_{r \in \F_q} (-1)^{\tr(ar)} \, \hat{W}^r(u)\,,
	\]
	and define the projector
	\begin{equation}\label{eq:qld-point-obs-def}
		\tau^{W,u}_a \,=\, \tau^W_{[g_\cdot(u) = a]}
		\,=\, \sum_{h \in \F_q^{M} \colon h \cdot \mathrm{ind}_m(u) = a} \tau^W_h\,,
	\end{equation}
	where $\{\tau^W_h\}_{h \in \F_q^{M}}$ are the joint eigenprojectors of
	the generalized Pauli $W$-observables
	(see Lemma~\ref{lem:pauli-observable-expansion}).
\end{definition}

\begin{definition}[Ancillary point measurement]
	\label{def:ancillary-point-measurement}
	\lean{MIPStarRE.QPBT.ProjectiveSetting.tauPointMeas}
	\lean{MIPStarRE.QPBT.ProjectiveSetting.tauPointMeas_effect}
	\leanok
	\uses{def:expanded-point-measurement}
	The ancillary point measurement at $W\in\{X,Z\}$ and $u\in\F_q^m$
	is the POVM $\{\tau^{W,u}_{a''}\}_{a''\in\F_q}$ on one Pauli
	register. It is obtained by measuring the generalized Pauli basis
	and reporting $a''=h\cdot\mathrm{ind}_m(u)$. Its effects are
	positive and sum to the identity, since the fibres of this map
	partition the Pauli basis labels.
\end{definition}

\begin{definition}[Product of strategy and ancillary point measurements]
	\label{def:point-tau-measurement}
	\lean{MIPStarRE.QPBT.ProjectiveSetting.pointTauMeas}
	\lean{MIPStarRE.QPBT.ProjectiveSetting.pointTauMeas_effect}
	\leanok
	\uses{def:ancillary-point-measurement, def:strategy-observables}
	For either player, the product measurement on the strategy register
	and one Pauli register has separate outcomes $(a',a'')\in\F_q^2$
	and effects
	\[
	N^{W,u}_{a',a''}
	=M^{(\mathrm{Point},W),u}_{a'}\otimes\tau^{W,u}_{a''}.
	\]
	This auxiliary measurement records the two outcomes whose sum defines
	the expanded point outcome. Tensor products preserve positivity, and
	completeness of both factors gives $\sum_{a',a''}N^{W,u}_{a',a''}=\Id$.
\end{definition}

\begin{lemma}[Convolution and projectivity of expanded point measurements]
	\label{lem:expanded-point-measurement-properties}
	\lean{MIPStarRE.QPBT.ProjectiveSetting.expPointOp_eq_convolution}
	\lean{MIPStarRE.QPBT.ProjectiveSetting.pointMeasExp_effect_eq_pointTauMeas_postprocess}
	\lean{MIPStarRE.QPBT.ProjectiveSetting.expPointOp_nonneg}
	\lean{MIPStarRE.QPBT.ProjectiveSetting.expPointOp_sum_eq_one}
	\lean{MIPStarRE.QPBT.ProjectiveSetting.pointMeasExp_isProjective}
	\leanok
	\uses{def:expanded-point-measurement, def:point-tau-measurement}
	For either player, $W \in \{X,Z\}$, $u \in \F_q^m$ and $a \in \F_q$,
	the expanded point effects have the convolution form
	\[
		\hat{M}^{(\mathrm{Point},W),u}_a
		\,=\, \sum_{\substack{a',a'' \in \F_q \\ a' + a'' = a}}
		M^{(\mathrm{Point},W),u}_{a'} \ot \tau^{W,u}_{a''}\,.
	\]
	Equivalently, the expanded point measurement is the postprocessing
	of $\{N^{W,u}_{a',a''}\}$ by the addition map $(a',a'')\mapsto a'+a''$.
	Each effect is positive semidefinite, and
	$\sum_{a \in \F_q}\hat{M}^{(\mathrm{Point},W),u}_a=\Id$.
	Moreover, each $\hat{M}^{(\mathrm{Point},W),u}_a$ is a projection
	and $\{\hat{M}^{(\mathrm{Point},W),u}_a\}_{a \in \F_q}$ is a projective
	measurement.
\end{lemma}
\begin{proof}
	\leanok
	\uses{def:strategy-observables, def:expanded-observables,
	lem:pauli-observable-expansion, prop:fourier-fact-scalar,
	lem:operator-average-scalar-constant, lem:point-observable-multiplication,
	lem:twisted-commutation, def:point-tau-measurement}
	Expand both tensor factors of $\hat{W}^r(u)$ and average over $r$.
	Character orthogonality retains exactly the terms with $a'+a''=a$,
	giving the convolution formula and the addition-postprocessing identity.
	Its summands are tensor products of positive semidefinite operators,
	proving positivity. Summing over $a$
	removes the convolution constraint; completeness of the strategy and
	Pauli measurements then gives the identity.
	For projectivity, multiplication of point observables and of Pauli
	observables gives $\hat{W}^r(u)\hat{W}^s(u)=\hat{W}^{r+s}(u)$.
	Squaring the Fourier average defining an effect and translating the
	uniform summation variable therefore gives the same effect.
	Positivity supplies self-adjointness, so each effect is a projection.
\end{proof}

\begin{definition}[Trace-coarse-grained point projections]
	\label{def:expanded-point-trace-projection}
	\lean{MIPStarRE.QPBT.ProjectiveSetting.expPointTrace}
	\leanok
	\uses{def:expanded-point-measurement, lem:expanded-point-measurement-properties,
	def:expanded-observables, def:subfield-trace}
	For fixed $u \in \F_q^m$, $r \in \F_q$ and $b \in \F_2$, define the
	projection
	\begin{equation}\label{eq:qld-def-mptur}
		\hat{M}^{(\mathrm{Point},W),u,r}_b \,=\, \sum_{a \in \F_q \colon \tr(ar) = b} \hat{M}^{(\mathrm{Point},W),u}_a\,,
	\end{equation}
	The register conventions are those of
	Definition~\ref{def:expanded-observables}.
\end{definition}

\begin{lemma}[Trace coarse-graining identity]
	\label{lem:expanded-point-trace-half-add}
	\lean{MIPStarRE.QPBT.ProjectiveSetting.expPointTrace_eq_half_add}
	\leanok
	\uses{def:expanded-point-trace-projection, def:expanded-observables}
	For either player, $W \in \{X,Z\}$, $u \in \F_q^m$, every $r \in \F_q$
	(including $r=0$), and $b \in \F_2$, we have
	\begin{equation}\label{eq:lc-22}
		\hat{M}^{(\mathrm{Point},W),u,r}_b \,=\, \frac{1}{2} \Big( \Id + (-1)^b\, \hat{W}^r(u) \Big)\,.
	\end{equation}
\end{lemma}
\begin{proof}
	\leanok
	\uses{lem:expanded-point-measurement-properties, prop:fourier-fact-scalar,
	lem:operator-average-scalar-constant}
	Fourier inversion gives
	$\sum_a(-1)^{\tr(ar)}\hat{M}^{(\mathrm{Point},W),u}_a=\hat{W}^r(u)$.
	The two trace classes sum to $\Id$ by completeness, and their difference
	is this signed sum. Adding or subtracting these two identities and
	dividing by two proves the formula.
\end{proof}

\begin{definition}[Symmetric equivalents]
	\label{def:symmetric-equivalents}
	\lean{MIPStarRE.QPBT.Placement, MIPStarRE.QPBT.Placement.IsOpposite}
	\lean{MIPStarRE.QPBT.ProjectiveSetting.place}
	\leanok
	\uses{def:expanded-state}
	For the remainder of this chapter, as well as in
	Chapters~\ref{chap:qpbt-combining} and~\ref{chap:qpbt-extraction}, the
	six registers
	$\mathsf{A}\mathsf{A}'\mathsf{A}''\mathsf{B}\mathsf{B}'\mathsf{B}''$
	are partitioned into two ``parties'' in two different ways. The first
	scheme takes $\mathsf{A}\mathsf{A}'$ as one party and
	$\mathsf{B}\mathsf{A}''$ as the other; the second scheme takes
	$\mathsf{B}\mathsf{B}'$ as the first party and
	$\mathsf{A}\mathsf{B}''$ as the second. The \emph{symmetric
	equivalents} of a bipartite relation of the form
	$M_{\mathsf{A}\mathsf{A}'} \approx_\delta N_{\mathsf{B}\mathsf{A}''}$
	are the three relations
	\[
		N_{\mathsf{A}\mathsf{A}'} \approx_\delta M_{\mathsf{B}\mathsf{A}''}\,,\qquad
		M_{\mathsf{B}\mathsf{B}'} \approx_\delta N_{\mathsf{A}\mathsf{B}''}\,,\qquad
		N_{\mathsf{B}\mathsf{B}'} \approx_\delta M_{\mathsf{A}\mathsf{B}''}\,,
	\]
	obtained by passing to the second bipartitioning scheme and by
	interchanging the first and second party; the symmetric equivalents of
	a single-party relation of the form
	$M_{\mathsf{A}\mathsf{A}'} \approx N_{\mathsf{A}\mathsf{A}'}$ are the
	same relation on the register pairs $\mathsf{B}\mathsf{A}''$,
	$\mathsf{B}\mathsf{B}'$ and $\mathsf{A}\mathsf{B}''$.
\end{definition}

\begin{lemma}[Transfer to symmetric equivalents]
	\label{lem:symmetric-equivalents-transfer}
	\notready
	\uses{def:symmetric-equivalents, def:expanded-state, lem:projective-strategy-setup, def:symmetric-game, def:EPR, def:povm-distance}
	Let $\mathcal{S} = (\ket{\psi}, M)$ be the projective strategy under
	analysis (Lemma~\ref{lem:projective-strategy-setup} and the convention
	following it), assumed symmetric in the sense of
	Definition~\ref{def:symmetric-game}: the state $\ket{\psi}$ is
	invariant under exchanging registers $\mathsf{A}$ and $\mathsf{B}$, and
	both players measure the same operators. Let $\ket{\hat{\psi}}$ be the
	expanded state~\eqref{eq:def-psihat}, let
	\[
		\sigma \colon\;
		\mathsf{A} \leftrightarrow \mathsf{B},\quad
		\mathsf{A}' \leftrightarrow \mathsf{B}',\quad
		\mathsf{A}'' \leftrightarrow \mathsf{B}''\,,
		\qquad\qquad
		\theta \colon\;
		\mathsf{A} \leftrightarrow \mathsf{B},\quad
		\mathsf{A}' \leftrightarrow \mathsf{A}''\,,
	\]
	be the two permutations of the six registers
	$\mathsf{A}\mathsf{A}'\mathsf{A}''\mathsf{B}\mathsf{B}'\mathsf{B}''$
	given by these transpositions, and let $U_\sigma$ and $U_\theta$ be the
	corresponding permutation unitaries. Write
	$\Pi = \{\Id,\, U_\sigma,\, U_\theta,\, U_\sigma U_\theta\}$ and let
	\[
		\mathcal{R} = \{\mathsf{A}\mathsf{A}',\; \mathsf{B}\mathsf{A}'',\;
		\mathsf{B}\mathsf{B}',\; \mathsf{A}\mathsf{B}''\}
	\]
	be the four register placements of
	Definition~\ref{def:symmetric-equivalents}. Then:
	\begin{enumerate}
		\item Every $U \in \Pi$ fixes $\ket{\hat{\psi}}$, and the induced
		action of $\Pi$ on $\mathcal{R}$ is simply transitive.
		\item Let $P$ and $Q$ be operators on
		$\hilb \ot (\C^q)^{\ot M}$, where $\hilb$ is the common Hilbert
		space of registers $\mathsf{A}$ and $\mathsf{B}$, let
		$\mathsf{C}\mathsf{C}', \mathsf{D}\mathsf{D}' \in \mathcal{R}$ and
		let $U \in \Pi$. Then
		\[
			\bigl\| \bigl( (P)_{\mathsf{C}\mathsf{C}'} - (Q)_{\mathsf{D}\mathsf{D}'} \bigr) \ket{\hat{\psi}} \bigr\|
			\,=\,
			\bigl\| \bigl( (P)_{U(\mathsf{C}\mathsf{C}')} - (Q)_{U(\mathsf{D}\mathsf{D}')} \bigr) \ket{\hat{\psi}} \bigr\|\,,
		\]
		the placement notation being that of
		Definition~\ref{def:expanded-observables}. Consequently, for
		families $\{P^x_a\}$ and $\{Q^x_a\}$ indexed by a question $x$ drawn
		from a distribution and an answer $a$, and for every $\delta \ge 0$,
		$(P^x_a)_{\mathsf{C}\mathsf{C}'} \approx_\delta (Q^x_a)_{\mathsf{D}\mathsf{D}'}$
		holds on $\ket{\hat{\psi}}$ if and only if
		$(P^x_a)_{U(\mathsf{C}\mathsf{C}')} \approx_\delta (Q^x_a)_{U(\mathsf{D}\mathsf{D}')}$
		does, with the same $\delta$
		(Definition~\ref{def:povm-distance}).
	\end{enumerate}
	In particular, taking $\mathsf{C}\mathsf{C}' = \mathsf{A}\mathsf{A}'$
	and $\mathsf{D}\mathsf{D}' = \mathsf{B}\mathsf{A}''$, respectively
	$\mathsf{D}\mathsf{D}' = \mathsf{A}\mathsf{A}'$, the images under
	$U_\theta$, $U_\sigma$ and $U_\sigma U_\theta$ are exactly the three
	symmetric equivalents of Definition~\ref{def:symmetric-equivalents},
	with the same error parameter.
\end{lemma}
\begin{proof}
	\uses{def:expanded-state, def:EPR, def:symmetric-game, def:povm-distance}
	Each of $\sigma$ and $\theta$ transposes registers carrying isomorphic
	Hilbert spaces --- $\mathsf{A}$ and $\mathsf{B}$ both carry $\hilb$,
	and the four ancilla registers are all isomorphic to
	$(\C^q)^{\ot M}$ (Definition~\ref{def:expanded-state}) --- so the
	permutation unitaries $U_\sigma$ and $U_\theta$ are defined and satisfy
	$U (P)_{\mathsf{C}\mathsf{C}'} U^\dagger = (P)_{U(\mathsf{C})U(\mathsf{C}')}$
	for every operator $P$ and every placement.

	For item~1, both generators fix $\ket{\hat{\psi}} = (\ket{\psi})_{\mathsf{A}\mathsf{B}} \ot \ket{\mathrm{EPR}_q}^{\ot M}_{\mathsf{A}'\mathsf{A}''} \ot \ket{\mathrm{EPR}_q}^{\ot M}_{\mathsf{B}'\mathsf{B}''}$:
	the factor $(\ket{\psi})_{\mathsf{A}\mathsf{B}}$ is invariant under
	$\mathsf{A} \leftrightarrow \mathsf{B}$ by the symmetry hypothesis;
	$\sigma$ exchanges the two ancilla factors, which are identical; and
	$\theta$ leaves $\ket{\mathrm{EPR}_q}^{\ot M}_{\mathsf{B}'\mathsf{B}''}$
	untouched and exchanges the two halves of each pair in
	$\ket{\mathrm{EPR}_q}^{\ot M}_{\mathsf{A}'\mathsf{A}''}$, under which
	$\ket{\mathrm{EPR}_q}$ is invariant (Definition~\ref{def:EPR}). The
	action on $\mathcal{R}$ sends $\mathsf{A}\mathsf{A}'$ to
	$\mathsf{A}\mathsf{A}'$, $\mathsf{B}\mathsf{B}'$,
	$\mathsf{B}\mathsf{A}''$ and $\mathsf{A}\mathsf{B}''$ under $\Id$,
	$U_\sigma$, $U_\theta$ and $U_\sigma U_\theta$ respectively, so it is
	simply transitive on the four placements.

	For item~2, insert $U^\dagger U$ and use $U \ket{\hat{\psi}} = \ket{\hat{\psi}}$
	together with unitary invariance of the norm:
	\[
		\bigl\| \bigl( (P)_{\mathsf{C}\mathsf{C}'} - (Q)_{\mathsf{D}\mathsf{D}'} \bigr) \ket{\hat{\psi}} \bigr\|
		= \bigl\| U \bigl( (P)_{\mathsf{C}\mathsf{C}'} - (Q)_{\mathsf{D}\mathsf{D}'} \bigr) U^\dagger \, U \ket{\hat{\psi}} \bigr\|
		= \bigl\| \bigl( (P)_{U(\mathsf{C}\mathsf{C}')} - (Q)_{U(\mathsf{D}\mathsf{D}')} \bigr) \ket{\hat{\psi}} \bigr\|\,.
	\]
	Summing the squares over the answers $a$ and averaging over the
	questions $x$ preserves the identity, so the two sides of the asserted
	equivalence of $\approx_\delta$ relations
	(Definition~\ref{def:povm-distance}) are equal quantities. The final
	assertion is item~1 applied to the placements listed there.
\end{proof}

Several lemmas below conclude by asserting that all symmetric equivalents
of the stated approximations also hold, in each case by
Lemma~\ref{lem:symmetric-equivalents-transfer} unless the individual register
pairs are established separately, as in Lemmas~\ref{lem:qld-comm-cons}
and~\ref{lem:qld-comm-line-cons}; later chapters invoke exactly these
variants.

\begin{lemma}[Consistency and approximate commutation of the expanded point measurements]
	\label{lem:qld-comm-cons}
	\lean{MIPStarRE.QPBT.deltaAnticom, MIPStarRE.QPBT.deltaAnticom_isPolyErr}
	\lean{MIPStarRE.QPBT.expPoint_self_cons, MIPStarRE.QPBT.expPointTrace_comm}
	\lean{MIPStarRE.QPBT.exists_deltaAnticom}
	\leanok
	\uses{def:expanded-point-measurement, def:expanded-point-trace-projection, def:expanded-state, def:approx-question-indexed-operators, def:povm-distance, def:symmetric-equivalents}
	Here $\poly(\eps)$ uses the corrected one-parameter convention: there
	exist $A \geq 1$ and $r>0$ such that
	$0 \leq \delta(x) \leq A x^r$ for every $x \geq 0$. Separating the
	prefactor and exponent is necessary to include the square-root bound;
	see \cite{gap:qpbt_polynomial-error-square-root}.
	There exists a function $\delta_{\mathrm{Anticom}}(\eps) = \poly(\eps)$
	such that the following holds, with respect to the state
	$\ket{\hat{\psi}}$ of Definition~\ref{def:expanded-state}.
	\begin{enumerate}
		\item (\textbf{Self-consistency}) For all $W \in \{X,Z\}$, on
		average over $r \in \F_q$ and $u \in \F_q^m$ (the relation does not
		depend on $r$), we have
		\[
			(\hat{M}^{(\mathrm{Point},W),u}_a)_{\mathsf{A}\mathsf{A}'}
			\,\approx_\eps\,
			(\hat{M}^{(\mathrm{Point},W),u}_a)_{\mathsf{B}\mathsf{A}''}\,,
		\]
		where the answer summation is over $a \in \F_q$.
		\item (\textbf{Approximate commutation}) For all $b, b' \in \F_2$,
		on average over uniformly random
		$\omega = (u_X,u_Z,r_X,r_Z) \in (\F_q^m)^2 \times (\F_q)^2$, we
		have
		\[
			(\hat{M}^{(\mathrm{Point},X),u_X,r_X}_{b}\,
			\hat{M}^{(\mathrm{Point},Z),u_Z,r_Z}_{b'})_{\mathsf{A}\mathsf{A}'}
			\,\approx_{\delta_{\mathrm{Anticom}}}\,
			(\hat{M}^{(\mathrm{Point},Z),u_Z,r_Z}_{b'}\,
			\hat{M}^{(\mathrm{Point},X),u_X,r_X}_{b})_{\mathsf{A}\mathsf{A}'}\,.
		\]
	\end{enumerate}
	Furthermore, all symmetric equivalents of these approximations also
	hold, and the proof yields $\delta_{\mathrm{Anticom}}(\eps) = \sqrt{\eps}$.
\end{lemma}
\begin{proof}
	\leanok \uses{lem:qld-win-implications, lem:qld-win-implications-obs, fact:data-processing, lem:twisted-commutation, def:expanded-point-measurement, def:expanded-point-trace-projection, def:expanded-observables, def:symmetric-equivalents,
	lem:expanded-point-trace-half-add, lem:qld-placed-consistency-distance,
	lem:expanded-point-measurement-properties}
	For item~1: the original point measurements
	$\{M^{(\mathrm{Point},W),u}_a\}$ are $\eps$-self-consistent between
	registers $\mathsf{A}$ and $\mathsf{B}$ in $\ket{\hat{\psi}}$
	(item~\ref{enu:qld-cons} of Lemma~\ref{lem:qld-win-implications}), the
	Pauli projections $\{\tau^{W,u}_a\}$ are perfectly self-consistent
	between registers $\mathsf{A}'$ and $\mathsf{A}''$ on the EPR ancillas,
	and the data processing inequality (Lemma~\ref{fact:data-processing})
	bounds the consistency defect of their addition postprocessing.
	By Lemma~\ref{lem:expanded-point-measurement-properties}, this
	postprocessing is the expanded point measurement in convolution form;
	Lemma~\ref{lem:qld-placed-consistency-distance} converts its consistency
	bound to the required state-dependent distance bound. For item~2:
	by Lemma~\ref{lem:expanded-point-trace-half-add}
	(\eqref{eq:lc-22}), each side expands as
	$\frac14 \big( \Id + (-1)^b \hat{X}^{r_X}(u_X) + (-1)^{b'} \hat{Z}^{r_Z}(u_Z) + (-1)^{b+b'} \hat{X}^{r_X}(u_X) \hat{Z}^{r_Z}(u_Z) \big)$
	respectively its reverse-ordered analogue, so it suffices to commute
	the two expanded observables. By~\eqref{eq:lc-23},
	$\hat{X}^{r_X}(u_X)\hat{Z}^{r_Z}(u_Z) = X^{r_X}(u_X) Z^{r_Z}(u_Z) \ot \tau^X(r_X \cdot \mathrm{ind}_m(u_X)) \tau^Z(r_Z \cdot \mathrm{ind}_m(u_Z))$;
	the strategy observables approximately commute or anticommute with sign
	$(-1)^{\gamma(\omega)}$ and error $\sqrt{\eps}$
	(\eqref{eq:pts-obs-commutation} of
	Lemma~\ref{lem:qld-win-implications-obs}), while the generalized Pauli
	observables exactly commute or anticommute with the same sign
	(see Lemma~\ref{lem:twisted-commutation}). The two signs cancel, so
	$\hat{X}^{r_X}(u_X)\hat{Z}^{r_Z}(u_Z) \approx_{\sqrt{\eps}} \hat{Z}^{r_Z}(u_Z)\hat{X}^{r_X}(u_X)$
	on average over uniformly random $\omega$, giving the approximate
	commutation with $\delta_{\mathrm{Anticom}} = \sqrt{\eps}$. The
	symmetric equivalents are obtained without the transfer of
	Lemma~\ref{lem:symmetric-equivalents-transfer}: both displays are proved
	separately on each of the four register pairs
	$\mathsf{A}\mathsf{A}'$, $\mathsf{B}\mathsf{A}''$, $\mathsf{B}\mathsf{B}'$ and
	$\mathsf{A}\mathsf{B}''$ of Definition~\ref{def:symmetric-equivalents}, so no
	symmetry hypothesis on the strategy is used.
\end{proof}

\begin{lemma}[Restriction of the low-degree encoding to a line]
	\label{lem:low-degree-encoding-line-restriction}
	\lean{MIPStarRE.QPBT.polynomialOnLine_lowDegreeEncoding_natDegree_le_sum}
	\lean{MIPStarRE.QPBT.polynomialOnLine_lowDegreeEncoding_natDegree_le}
	\lean{MIPStarRE.QPBT.polynomialOnLine_lowDegreeEncoding_natDegree_le_one_of_axis}
	\lean{MIPStarRE.QPBT.restrictToLine_lowDegreeEncoding_fitsDegree_of_axis}
	\lean{MIPStarRE.QPBT.DegPoly.FitsDegree}
	\lean{MIPStarRE.QPBT.eval_polynomialOnLine, MIPStarRE.QPBT.evalCoefficient_restrictToLine}
	\lean{MIPStarRE.QPBT.evaluatesTo_restrictToLine, MIPStarRE.QPBT.evalOpt_restrictToLine}
	\lean{MIPStarRE.QPBT.evalOpt_restrictToLine_lowDegreeEncoding}
	\lean{MIPStarRE.QPBT.restrictToAxisLine}
	\lean{MIPStarRE.QPBT.evalCoefficient_restrictToAxisLine_lowDegreeEncoding}
	\leanok
	\uses{def:low-degree-encoding, def:indicator-vector, def:ideg-deg-polynomials, def:line}
	Let $g_h$ be the multilinear low-degree encoding of
	$h\in\F_q^{\{0,1\}^m}$, and let $\ell$ have base point $u_0$ and
	direction $v$. The univariate polynomial
	$f_\ell(t)=g_h(u_0+tv)$ has degree at most $md$.
	If $\ell$ is axis-parallel, its degree is at most $1\leq d$; hence its
	coefficient list, regarded as an element of $\deg_{md}(\ell)$, has
	zero entries in every degree greater than $d$.
	Accordingly, on an axis-parallel line the restriction may equally be
	recorded as a coefficient list in $\deg_{d}(\ell)$, and this degree-$d$
	list has the same coefficient-list evaluation $g_h(u)$ at every
	$u\in\ell$; it is the form reported for $\mathsf{ALine}$ answers in the
	proof of Lemma~\ref{lem:pauli-completeness}.
	At every point $u\in\ell$, its coefficient-list evaluation is
	$g_h(u)=h\cdot\mathrm{ind}_m(u)$. This assertion also applies when
	$v=0$: every parameter represents the same point and gives the same
	value.
	More generally, if the affine polynomial $u_{0,i}+tv_i$ has degree
	at most $b_i$ for each coordinate $i$, then $f_\ell$ has degree at
	most $\sum_i b_i$.

	This auxiliary records the restriction calculation used for honest line
	answers in the proof of Lemma~\ref{lem:pauli-completeness} and for the
	ancillary line measurement below. For a general multivariate polynomial
	$g$, substitution still gives
	$f_\ell(t)=g(u_0+tv)$; passage to its degree-$md$ coefficient list
	preserves evaluation provided the substituted polynomial has degree at
	most $md$.
\end{lemma}

\begin{proof}
	\leanok
	Each factor in an indicator polynomial from
	Definition~\ref{def:low-degree-encoding} becomes a univariate polynomial
	of degree at most one after affine substitution. Each product therefore
	has degree at most $m$, and a linear combination has the same bound.
	On an axis-parallel line only one factor can be nonconstant, giving the
	degree bound one. The degree bounds imply that the omitted coefficients
	vanish, so coefficient-list evaluation is ordinary polynomial
	evaluation. Substitution then gives $g_h(u)$, which equals
	$h\cdot\mathrm{ind}_m(u)$ by
	Equation~\eqref{eq:low-degree-encoding-definition}.
\end{proof}

\begin{definition}[Expanded line measurements]
	\label{def:expanded-line-measurement}
	\lean{MIPStarRE.QPBT.polynomialOnLine, MIPStarRE.QPBT.restrictToLine}
	\lean{MIPStarRE.QPBT.polynomialOnLine_lowDegreeEncoding_natDegree_le}
	\lean{MIPStarRE.QPBT.tauLineProj, MIPStarRE.QPBT.ProjectiveSetting.expLineOp}
	\lean{MIPStarRE.QPBT.ProjectiveSetting.lineMeasExp}
	\lean{MIPStarRE.QPBT.ProjectiveSetting.lineMeasExp_isProjective}
	\lean{MIPStarRE.QPBT.ProjectiveSetting.expLineOp_zero_of_not_deg_d}
	\lean{MIPStarRE.QPBT.tauLineMeas}
	\leanok
	% In the source this measurement is defined inside the proof of the lemma
	% labelled lem:qld-comm-line-cons (lines 552-556); it is recorded as a
	% standalone definition because Chapter chap:qpbt-combining uses the
	% concrete measurement, not merely its existence.
	\uses{lem:projective-strategy-setup, def:expanded-state, def:expanded-observables, def:ideg-deg-polynomials, def:low-degree-encoding, lem:pauli-observable-expansion, def:line, def:pauli-win-predicate}
	For $W \in \{X,Z\}$, a line $\ell$, and $f \in \deg_{md}(\ell)$,
	define
	\[
		\hat{M}^{(\mathrm{Line},W),\ell}_f
		\,=\, \sum_{\substack{f',f'' \in \deg_{md}(\ell) \\ f' + f'' = f}}
		M^{(\mathrm{Line},W),\ell}_{f'} \ot \tau^{W,\ell}_{f''}\,,
	\]
	where the projector $\tau^{W,\ell}_{f''}$, acting on register
	$\mathsf{A}'$ (or on $\mathsf{A}''$, following the register conventions
	of Definition~\ref{def:expanded-observables}), is the outcome-$f''$
	projector of the following projective measurement: measure all $M$
	qudits in the basis $W$, obtaining an outcome $h \in \F_q^{M}$, and
	return the polynomial $f''$ obtained by substituting the
	parameterization of $\ell$ into the low-degree encoding
	$g_h \colon \F_q^m \to \F_q$
	(see Definition~\ref{def:low-degree-encoding}): writing $(u_0,v)$ for
	the base point and direction of $\ell$, the outcome $f''$ is
	$g_h(u_0 + tv)$, expanded as a polynomial in the variable $t$ from the
	coefficient representation of $g_h$. Since $g_h$ has individual degree
	at most $d$, hence total degree at most $md$, the substituted
	polynomial has degree at most $md$ in $t$, so
	$f'' \in \deg_{md}(\ell)$. The strategy measurement
	$\{M^{(\mathrm{Line},W),\ell}_{f'}\}$ is indexed by the answers of the
	prescribed form (Definition~\ref{def:pauli-win-predicate}): for a
	diagonal line these are the elements of $\deg_{md}(\ell)$, and for an
	axis-parallel line the elements of $\deg_d(\ell)$, regarded as a
	measurement with outcomes in $\deg_{md}(\ell)$ through the inclusion
	$\deg_d(\ell) \subseteq \deg_{md}(\ell)$ of
	Definition~\ref{def:ideg-deg-polynomials}, with the zero operator
	assigned to outcomes outside $\deg_d(\ell)$. For an axis-parallel
	$\ell$, with direction $e_i$, the substituted polynomial also has
	degree at most $d$ in $t$ --- its degree is the individual degree of
	$g_h$ in the $i$-th variable --- so
	$\hat{M}^{(\mathrm{Line},W),\ell}_f = 0$ unless $f \in \deg_d(\ell)$.
	In all cases
	$\{\hat{M}^{(\mathrm{Line},W),\ell}_f\}_{f \in \deg_{md}(\ell)}$ is a
	projective measurement.
\end{definition}

\begin{definition}[Auxiliary strategy--Pauli product measurement on a line]
	\label{def:line-tau-product-measurement}
	\lean{MIPStarRE.QPBT.ProjectiveSetting.lineTauMeas}
	\lean{MIPStarRE.QPBT.ProjectiveSetting.lineTauMeas_effect}
	\leanok
	\uses{def:expanded-line-measurement, def:ideg-deg-polynomials}
	For $W \in \{X,Z\}$ and a line $\ell$, the product measurement
	\[
		\hat{N}^{(\mathrm{Line},W),\ell}_{(f',f'')}
		\,=\, M^{(\mathrm{Line},W),\ell}_{f'} \ot \tau^{W,\ell}_{f''}\,,
		\qquad (f',f'') \in \deg_{md}(\ell) \times \deg_{md}(\ell),
	\]
	has outcome set $\deg_{md}(\ell) \times \deg_{md}(\ell)$ and retains the
	strategy polynomial and the Pauli line polynomial as separate outcomes;
	completeness follows from completeness of the two factors.
	This auxiliary measurement is used in the proof of
	Lemma~\ref{lem:qld-comm-line-cons}; its addition postprocessing is identified
	in Lemma~\ref{lem:expanded-line-measurement-postprocess}.
\end{definition}

\begin{lemma}[Addition postprocessing of the line product measurement]
	\label{lem:expanded-line-measurement-postprocess}
	\lean{MIPStarRE.QPBT.ProjectiveSetting.lineMeasExp_effect_eq_lineTauMeas_postprocess}
	\leanok
	\uses{def:expanded-line-measurement, def:line-tau-product-measurement}
	For either player, $W \in \{X,Z\}$ and a line $\ell$, the image of the
	product measurement of Definition~\ref{def:line-tau-product-measurement}
	under the addition map $(f',f'') \mapsto f'+f''$ is the expanded line
	measurement. Thus, for every $f \in \deg_{md}(\ell)$,
	\[
		\hat{M}^{(\mathrm{Line},W),\ell}_f
		\,=\, \sum_{f'+f''=f} \hat{N}^{(\mathrm{Line},W),\ell}_{(f',f'')}\,.
	\]
\end{lemma}
\begin{proof}
	\leanok
	\uses{def:expanded-line-measurement, def:line-tau-product-measurement}
	Expand the effect of the postprocessed measurement. Its fibre over $f$
	consists of the pairs $(f',f'')$ satisfying $f'+f''=f$, so the resulting sum
	is exactly the convolution defining the expanded line effect.
\end{proof}

\begin{lemma}[Consistency of the expanded line measurements]
	\label{lem:qld-comm-line-cons}
	\lean{MIPStarRE.QPBT.deltaLine, MIPStarRE.QPBT.deltaLine_isPolyErr}
	\lean{MIPStarRE.QPBT.expLine_self_cons, MIPStarRE.QPBT.expLine_point_cons}
	\lean{MIPStarRE.QPBT.expLine_point_cons', MIPStarRE.QPBT.exists_deltaLine}
	\lean{MIPStarRE.QPBT.ExpandedLineConclusions}
	\lean{MIPStarRE.QPBT.ProjectiveSetting.lineEvalMeasExp}
	\lean{MIPStarRE.QPBT.ProjectiveSetting.pointMeasExpOption}
	\lean{MIPStarRE.QPBT.ProjectiveSetting.expPointEffectAtLineAnswer}
	\leanok
	\uses{def:expanded-point-measurement, def:expanded-line-measurement, def:expanded-state, def:line-point-dist, def:approx-question-indexed-operators, def:povm-distance, def:ideg-deg-polynomials, def:symmetric-equivalents}
	Here $\poly(\eps)$ uses the corrected one-parameter convention: there
	exist $A \geq 1$ and $r>0$ such that
	$0 \leq \delta(x) \leq A x^r$ for every $x \geq 0$. Separating the
	prefactor and exponent is necessary to include the square-root bound;
	see \cite{gap:qpbt_polynomial-error-square-root}.
	There exists a function $\delta_{\mathrm{Line}}(\eps) = \poly(\eps)$
	such that the following holds, with respect to the state
	$\ket{\hat{\psi}}$ of Definition~\ref{def:expanded-state}. For all
	$W \in \{X,Z\}$ and lines $\ell$ there exists a projective measurement
	$\{\hat{M}^{(\mathrm{Line},W),\ell}_f\}_{f \in \deg_{md}(\ell)}$,
	acting on registers $\mathsf{A}\mathsf{A}'$ or
	$\mathsf{B}\mathsf{A}''$, such that:
	\begin{enumerate}
		\item \label{enu:qld-comm-line-self-cons}
		(\textbf{Self-consistency}) For all $W \in \{X,Z\}$ and
		$r \in \F_q$ (the relation does not depend on $r$), on average over
		a uniformly random pair $(\ell,u)$ drawn from the line-point
		distribution (Definition~\ref{def:line-point-dist}),
		\[
			(\hat{M}^{(\mathrm{Line},W),\ell}_f)_{\mathsf{A}\mathsf{A}'}
			\,\approx_{\delta_{\mathrm{Line}}}\,
			(\hat{M}^{(\mathrm{Line},W),\ell}_f)_{\mathsf{B}\mathsf{A}''}\,,
		\]
		where the answer summation is over $f \in \deg_{md}(\ell)$.
		\item \label{eq:qld-comm-line-pt-cons}
		(\textbf{Consistency with points measurements I}) On average over
		a uniformly random pair $(\ell,u)$ drawn from the line-point
		distribution,
		\[
			(\hat{M}^{(\mathrm{Line},W),\ell}_f)_{\mathsf{A}\mathsf{A}'}
			\,\approx_{\delta_{\mathrm{Line}}}\,
			(\hat{M}^{(\mathrm{Line},W),\ell}_f)_{\mathsf{A}\mathsf{A}'}
			\ot (\hat{M}^{(\mathrm{Point},W),u}_{f(u)})_{\mathsf{B}\mathsf{A}''}\,,
		\]
		where the answer summation is over $f \in \deg_{md}(\ell)$.
		\item \label{eq:qld-comm-line-pt-cons2}
		(\textbf{Consistency with points measurements II}) On average over
		a uniformly random pair $(\ell,u)$ drawn from the line-point
		distribution,
		\[
			(\hat{M}^{(\mathrm{Line},W),\ell}_{[\mathrm{eval}_u(\cdot)=a]})_{\mathsf{A}\mathsf{A}'}
			\,\approx_{\delta_{\mathrm{Line}}}\,
			(\hat{M}^{(\mathrm{Point},W),u}_a)_{\mathsf{B}\mathsf{A}''}\,,
		\]
		where, for $a \in \F_q$,
		$\hat{M}^{(\mathrm{Line},W),\ell}_{[\mathrm{eval}_u(\cdot)=a]} = \sum_{f \in \deg_{md}(\ell) \colon f(u)=a} \hat{M}^{(\mathrm{Line},W),\ell}_f$.
		For $a=\bot$ it is the sum over the polynomials without an evaluation
		at $u$, and the answer summation is over
		$a \in \F_q \cup \{\bot\}$.
	\end{enumerate}
	Furthermore, all symmetric equivalents of these approximations also
	hold. The proof exhibits the measurements of
	Definition~\ref{def:expanded-line-measurement} as witnesses, with the
	common error $\delta_{\mathrm{Line}}(\eps) = \sqrt{\eps}$ in all three
	items. The stronger linear-error assertion made in the source proof is
	stated separately in Lemma~\ref{lem:expanded-line-point-linear-error}.
\end{lemma}
\begin{proof}
	\leanok
	\uses{def:expanded-line-measurement, lem:expanded-line-measurement-postprocess,
	def:expanded-point-measurement, lem:qld-win-implications,
	lem:qld-reversed-consistency, fact:data-processing,
	fact:add-a-proj2, fact:agreement, def:line-point-dist,
	def:ideg-deg-polynomials, def:pauli-win-predicate, def:symmetric-equivalents}
	We first recall the source argument.
	Take $\{\hat{M}^{(\mathrm{Line},W),\ell}_f\}$ to be the measurement of
	Definition~\ref{def:expanded-line-measurement}.
	Item~\ref{enu:qld-comm-line-self-cons} follows from the fact that the
	original line measurements $\{M^{(\mathrm{Line},W),\ell}_f\}$ are
	$\eps$-self-consistent between registers $\mathsf{A}$ and $\mathsf{B}$
	in $\ket{\hat{\psi}}$ (item~\ref{enu:qld-cons} of
	Lemma~\ref{lem:qld-win-implications}), the Pauli projections
	$\{\tau^{W,\ell}_{f''}\}$ are perfectly self-consistent between
	registers $\mathsf{A}'$ and $\mathsf{A}''$, and the data processing
	inequality (Lemma~\ref{fact:data-processing}) applied to the addition
	postprocessing of Definition~\ref{def:line-tau-product-measurement}. For
	item~\ref{eq:qld-comm-line-pt-cons}: for the outcomes $f'$ that
	evaluate at $u$, insert the projection
	$M^{(\mathrm{Line},W),\ell}_{[\mathrm{eval}_u(\cdot)=f'(u)]}$ by
	projectivity, replace it by
	$(M^{(\mathrm{Point},W),u}_{f'(u)})_{\mathsf{B}\mathsf{A}''}$ using the
	line-point consistency (item~\ref{enu:qld-ld} of
	Lemma~\ref{lem:qld-win-implications}) together with
	Lemma~\ref{fact:add-a-proj2}, adjoin the exactly consistent projector
	$\tau^{W,u}_{f''(u)}$ on register $\mathsf{A}''$, and recombine by the
	convolution form of Definition~\ref{def:expanded-point-measurement}.
	Outcomes without an evaluation at $u$ occur only when $\ell$ has zero
	direction, where $\tau^{W,\ell}_{f''}$ vanishes unless $f''$ is a
	constant polynomial, so that $f = f' + f''$ evaluates at $u$ exactly
	when $f'$ does; for these outcomes the right-hand side of
	item~\ref{eq:qld-comm-line-pt-cons} carries the zero operator
	(Definition~\ref{def:ideg-deg-polynomials}), and their total mass
	satisfies
	\[
		\E_{(\ell,u)} \sum_{f' \text{ without an evaluation at } u}
		\bra{\hat{\psi}}\, M^{(\mathrm{Line},W),\ell}_{f'} \ot \Id\,
		\ket{\hat{\psi}} \,\leq\, O(\eps)\,,
	\]
	because the left-hand side is exactly the mass of the class
	$M^{(\mathrm{Line},W),\ell}_{[\mathrm{eval}_u(\cdot)=\bot]}$ of the
	completed family in item~\ref{enu:qld-ld} of
	Lemma~\ref{lem:qld-win-implications} --- its value on
	$\ket{\hat{\psi}}$ coincides with its value on $\ket{\psi}$, the
	operator acting trivially on the ancilla registers --- whose partner
	class $M^{(\mathrm{Point},W),u}_{\bot}$ is the zero operator, so that
	the whole of this mass is disagreement mass of that consistency
	statement.
	This gives item~\ref{eq:qld-comm-line-pt-cons} with error $\eps$. For
	item~\ref{eq:qld-comm-line-pt-cons2}, the data processing inequality
	does not apply directly, since the bound of
	item~\ref{eq:qld-comm-line-pt-cons} is stated with $\approx$ rather
	than $\simeq$ and its right-hand side is not of the form
	$\Id \ot B$; instead, a direct computation from
	item~\ref{eq:qld-comm-line-pt-cons} together with projectivity gives
	% The source ends this display with the exact upper bound eps; the bound
	% inherited from item 2 carries the hidden universal constant of
	% def:povm-distance, which is retained inside the O(.).
	\[
		\E_{(\ell,u)} \sum_{a \in \F_q}
		\bra{\hat{\psi}}\, \hat{M}^{(\mathrm{Line},W),\ell}_{[\mathrm{eval}_u(\cdot)=a]}
		\ot \big( \Id - \hat{M}^{(\mathrm{Point},W),u}_a \big) \ket{\hat{\psi}}
		\,\leq\, O(\eps)\,,
	\]
	while the mass bound of the previous paragraph controls the outcomes
	omitted from every class,
	$\E_{(\ell,u)} \bra{\hat{\psi}} \big( \Id - \sum_a
	\hat{M}^{(\mathrm{Line},W),\ell}_{[\mathrm{eval}_u(\cdot)=a]} \big)
	\ot \Id \ket{\hat{\psi}} \leq O(\eps)$; both terms are needed because
	for a zero-direction $\ell$ the classes
	$[\mathrm{eval}_u(\cdot)=a]$ need not exhaust $\deg_{md}(\ell)$.
	An elementary bound for projective sub-measurements $\{A_i\}$,
	$\{B_i\}$ and a state $\ket{\varphi}$, namely
	$\sum_i \bra{\varphi}(A_i - B_i)\ket{\varphi} \leq 2 \big( \sum_i \big\|(A_i - B_i)\ket{\varphi}\big\|^2 \big)^{1/2}$,
	upgrades these two bounds to
	$\Id \approx_{\sqrt{\eps}} \sum_a \hat{M}^{(\mathrm{Line},W),\ell}_{[\mathrm{eval}_u(\cdot)=a]} \ot \hat{M}^{(\mathrm{Point},W),u}_a$.
	A second elementary fact about projective measurements --- if
	$\Id \approx_\delta \sum_i A^x_i \ot B^x_i$ for families of projective
	measurements $\{A^x_i\}$ and $\{B^x_i\}$ acting on separate registers
	of a bipartite state, then $A^x_i \ot \Id \simeq_\delta \Id \ot B^x_i$,
	and hence $A^x_i \ot \Id \approx_\delta \Id \ot B^x_i$ by item~1 of
	Lemma~\ref{fact:agreement} --- is then applied to the two families
	completed as in Definition~\ref{def:ideg-deg-polynomials}: the
	evaluation classes of the expanded line measurement together with
	$\hat{M}^{(\mathrm{Line},W),\ell}_{[\mathrm{eval}_u(\cdot)=\bot]}$,
	and the expanded point measurement together with
	$\hat{M}^{(\mathrm{Point},W),u}_{\bot} = 0$. Both completed families
	are projective measurements, and the completion leaves the sum in the
	hypothesis unchanged, since the term for the outcome $\bot$ vanishes.
	Restricted to the outcomes in $\F_q$, the conclusion yields
	item~\ref{eq:qld-comm-line-pt-cons2} with error $\sqrt{\eps}$, so that
	$\delta_{\mathrm{Line}}(\eps) = \sqrt{\eps}$. The symmetric
	equivalents follow from
	Lemma~\ref{lem:symmetric-equivalents-transfer}.

	An equivalent argument reverses the order of the estimates. The exact
	Pauli overlap computation identifies the inconsistency of the expanded
	evaluation classes with that of the strategy's original completed
	evaluation classes. Applying the winning implication, its reversed form
	in Lemma~\ref{lem:qld-reversed-consistency}, and the
	agreement estimate gives item~\ref{eq:qld-comm-line-pt-cons2} with
	an error linear in $\eps$. Projective refinement then gives
	item~\ref{eq:qld-comm-line-pt-cons}, also with linear error.
	For each estimate, combining $x\leq a\eps$ with the universal bound
	$x\leq 4$ gives $x\leq(a+4)\sqrt{\eps}$: use
	$\eps\leq\sqrt{\eps}$ when $0\leq\eps\leq 1$ and
	$1\leq\sqrt{\eps}$ when $\eps\geq 1$.
	The four directed opposite-placement pairs follow separately from
	their explicit bipartition identities. Thus this argument does not
	rely on the symmetric-equivalent transfer lemma or the intermediate
	sub-measurement inequality in the preceding source argument.
\end{proof}

\begin{lemma}[Linear error in expanded line-point consistency]
	\label{lem:expanded-line-point-linear-error}
	\lean{MIPStarRE.QPBT.expLine_point_cons_linear}
	\leanok
	\uses{def:expanded-line-measurement, def:expanded-point-measurement,
	def:expanded-state, def:line-point-dist, def:povm-distance,
	def:ideg-deg-polynomials, def:symmetric-equivalents}
	Under the projective-strategy hypotheses of this section, the concrete
	measurements of Definition~\ref{def:expanded-line-measurement} satisfy
	\begin{equation}
		\label{eq:qld-comm-line-pt-cons-eps}
		(\hat{M}^{(\mathrm{Line},W),\ell}_f)_{\mathsf{A}\mathsf{A}'}
		\,\approx_{\eps}\,
		(\hat{M}^{(\mathrm{Line},W),\ell}_f)_{\mathsf{A}\mathsf{A}'}
		\ot (\hat{M}^{(\mathrm{Point},W),u}_{f(u)})_{\mathsf{B}\mathsf{A}''}
	\end{equation}
	on the line-point distribution, for every $W\in\{X,Z\}$ and all four
	directed opposite placements. The answer sum runs over
	$f\in\deg_{md}(\ell)$; the selected point effect is zero when $f$ has
	no evaluation at $u$. More precisely, the average squared distance is at
	most $C\eps$ for a universal $C\geq1$, independent of the parameters,
	strategy, basis and placements. This is the stronger assertion in the
	proof of the expanded-line consistency lemma of \cite{Ji2020MIPStar},
	% Source: eq:qld-comm-line-pt-cons-eps and the following paragraph.
	rather than the common-polynomial statement of
	Lemma~\ref{lem:qld-comm-line-cons}.
\end{lemma}
\begin{proof}
	\leanok
	\uses{lem:qld-win-implications, lem:qld-reversed-consistency,
	fact:agreement, fact:add-a-proj2, def:expanded-line-measurement,
	def:expanded-point-measurement, def:symmetric-equivalents}
	Let $D_{\mathrm I}$ be the distance in the displayed assertion and
	$D_{\mathrm{II}}$ the distance between the completed expanded
	evaluation classes and point measurements. Projective refinement gives
	$D_{\mathrm I}\leq D_{\mathrm{II}}$. The exact consistency of the
	ancillary Pauli projectors identifies their summed overlap with that
	of the original strategy measurements, including the $\bot$ class.
	The agreement inequality therefore gives
	$D_{\mathrm{II}}\leq 2\,\mathrm{Cons}_{\mathrm{strategy}}$.
	For a line on Alice's side the forward low-degree winning implication
	bounds this defect by $C_1\eps$; for a line on Bob's side use
	Lemma~\ref{lem:qld-reversed-consistency} to obtain $C_2\eps$.
	Both estimates apply on each ancillary bipartition. Taking
	$C=2(C_1+C_2)$ proves the assertion for all four directed pairs.
\end{proof}

The expanded point and line measurements constructed in this section, being
approximately consistent and approximately commuting, are combined into a
single strategy for the classical low individual degree test in
Chapter~\ref{chap:qpbt-combining}.

\section{Explicit quantitative observable estimates}
\label{sec:qld-observables-explicit}

This section records the numerical witnesses used by the current formal proof.
They are Lean-only quantitative refinements of the preceding source-facing
statements.  In particular, none of the asymptotic statements above is replaced
by one of the specializations below.

\begin{definition}[Observable constants of the quantitative baseline]
\label{def:pauli-observable-baseline-constants}
\lean{MIPStarRE.QPBT.pauliBaselineCommutatorConstant}
\lean{MIPStarRE.QPBT.pauliBaselineTwistedConstant}
\leanok
\uses{lem:pauli-edge-card-explicit}
Put
\[
 C_{\mathrm{comm}}=1024(172+\sqrt{344})
 \quad\text{and}\quad
 C_{\mathrm{tw}}=2(2C_{\mathrm{comm}}+4886363136+4).
\]
The factor two in $C_{\mathrm{tw}}$ is the loss incurred when the
phase-signed commutator is transferred to either expanded-space placement.
\end{definition}

\begin{theorem}[Explicit winning implications]
\label{thm:qld-win-implications-explicit}
\lean{MIPStarRE.QPBT.WinImplications.win_low_degree_explicit}
\lean{MIPStarRE.QPBT.WinImplications.win_low_degree_interchanged_explicit}
\lean{MIPStarRE.QPBT.WinImplications.win_pauli_basis_cons_explicit}
\lean{MIPStarRE.QPBT.WinImplications.win_pauli_basis_cons_interchanged_explicit}
\lean{MIPStarRE.QPBT.WinImplications.win_comm_explicit}
\lean{MIPStarRE.QPBT.WinImplications.win_comm_interchanged_explicit}
\lean{MIPStarRE.QPBT.WinImplications.win_comm_swapped_explicit}
\lean{MIPStarRE.QPBT.WinImplications.win_comm_cons_explicit}
\lean{MIPStarRE.QPBT.WinImplications.win_comm_cons_interchanged_explicit}
\lean{MIPStarRE.QPBT.WinImplications.win_comm_cons_swapped_explicit}
\lean{MIPStarRE.QPBT.WinImplications.win_magic_square_explicit}
\lean{MIPStarRE.QPBT.WinImplications.win_ms_cons_explicit}
\lean{MIPStarRE.QPBT.WinImplications.win_ms_cons_interchanged_explicit}
\lean{MIPStarRE.QPBT.WinImplications.win_ms_cons_swapped_explicit}
\leanok
\uses{lem:qld-win-implications, lem:pauli-edge-card-explicit}
Let $S$ be a projective strategy of value at least $1-\eps$, where
$\eps\geq0$.  The low-degree and point--Pauli consistency defects, in either
question orientation, are at most $86\eps$.  The commuting Pair consistency
defects and the commuting point-trace consistency defects are at most
$172\eps$ in the forward, interchanged, and swapped orientations.  The three
Magic Square variable-consistency defects are at most $1376\eps$, and
\[
 \left|1-\E_{\omega\sim\mathcal D_{\mathrm{anti}}}
   \operatorname{val}(S_{\omega}^{\mathrm{MS}})\right|
 \leq1376\eps.
\]
Each defect here is the corresponding defect displayed in
Lemma~\ref{lem:qld-win-implications}; the statement fixes its coefficient and
also records the two factor-interchanged forms used below.
\end{theorem}
\begin{proof}
\leanok
Each bad event is charged to one fixed verifier edge.  Conditioning on a
commuting tuple costs a factor $2$, conditioning on an anticommuting tuple
costs a factor $16$, and Lemma~\ref{lem:pauli-edge-card-explicit} supplies the
edge count $86$.
\end{proof}

\begin{theorem}[Explicit commutator and anticommutator estimates]
\label{thm:qld-point-observable-halves-explicit}
\lean{MIPStarRE.QPBT.WinImplications.point_trace_commutator_comm_le_explicit}
\lean{MIPStarRE.QPBT.WinImplications.point_obs_commutator_comm_le_explicit}
\lean{MIPStarRE.QPBT.WinImplications.point_obs_commutator_comm_le_alice_explicit}
\lean{MIPStarRE.QPBT.WinImplications.point_obs_commutator_comm_le_bob_explicit}
\lean{MIPStarRE.QPBT.WinImplications.point_obs_anticommutator_anticomm_le_explicit}
\lean{MIPStarRE.QPBT.WinImplications.point_obs_anticommutator_anticomm_le_bob_explicit}
\leanok
\uses{thm:qld-win-implications-explicit, fact:agreement}
Suppose two binary measurements are each within squared operator-family
distance $\delta$ of the corresponding marginal of one joint projective
measurement.  The average squared norm of their effect commutator is at
most $16\delta$.

More generally, let $\chi$ be a unit state, let $c_1\geq0$, let the joint
measurement be projective, and put $O_W=I-2M_W(1)$ for $W\in\{X,Z\}$.
If the two direct consistency defects are at
most $c_1$, the two detour defects are at most $c_2$, and the two defects from
the detour measurements to the joint marginals are at most $c_3$, then the
observable commutator has bound
\[
  1024\bigl(c_1+\sqrt{c_2+c_3}\bigr).
\]
For the concrete Pauli strategy, on Alice's state and on the swapped Bob
state, respectively, this gives
\[
 \E_{\mathcal D_{\mathrm{comm}}}\|[X,Z]\ket\psi\|^2
 \leq C_{\mathrm{comm}}(\eps+\sqrt\eps).
\]
On the anticommuting distribution, in the same two orientations,
\[
 \E_{\mathcal D_{\mathrm{anti}}}\|\{X,Z\}\ket\psi\|^2
 \leq4886363136\eps.
\]
\end{theorem}
\begin{proof}
\leanok
Apply the explicit joint-measurement commutator estimate, pass from
projections to reflections, and substitute the constants of
Theorem~\ref{thm:qld-win-implications-explicit}.  The anticommuting estimate
uses the robust Magic Square relation and the same consistency bounds.
\end{proof}

\begin{theorem}[Explicit phase-signed commutation]
\label{thm:qld-twisted-commutation-explicit}
\lean{MIPStarRE.QPBT.WinImplications.point_obs_twisted_commutation_explicit}
\lean{MIPStarRE.QPBT.WinImplications.point_obs_twisted_commutation_interchanged_explicit}
\lean{MIPStarRE.QPBT.ProjectiveSetting.twisted_commutator_avg_le_explicit}
\leanok
\uses{def:pauli-observable-baseline-constants,
thm:qld-point-observable-halves-explicit}
For either player, the average squared distance between $XZ$ and
$(-1)^{\gamma(\omega)}ZX$ is at most
\[
 (2C_{\mathrm{comm}}+4886363136+4)\sqrt\eps
\]
in the original strategy order, where
$\gamma(\omega)=\operatorname{tr}((r_X\operatorname{ind}_m(u_X))
\cdot(r_Z\operatorname{ind}_m(u_Z)))$ as in
Equation~\eqref{eq:pts-obs-commutation}.  After placement on any one of the
four expanded registers, the corresponding twisted-commutator average is at
most $C_{\mathrm{tw}}\sqrt\eps$.
\end{theorem}
\begin{proof}
\leanok
Split the uniform tuple distribution into its commuting and anticommuting
parts, use Theorem~\ref{thm:qld-point-observable-halves-explicit}, and use the
universal bound $4$ outside the range $\eps\leq1$.  Reindexing the swapped
state proves the second player orientation.
\end{proof}

\begin{theorem}[Explicit expanded point and line comparisons]
\label{thm:qld-expanded-comparisons-explicit}
\lean{MIPStarRE.QPBT.exp_point_self_cons_explicit}
\lean{MIPStarRE.QPBT.exp_line_point_cons_explicit'}
\leanok
\uses{lem:qld-comm-cons, lem:qld-comm-line-cons,
thm:qld-twisted-commutation-explicit}
For every basis and every directed pair of opposite placements, the expanded
point measurements have squared operator-family distance at most $172\eps$.
On average over the line--point distribution, the sum over completed answers
of the squared operator-action distances between the expanded line-evaluation
family and the opposite expanded point family is at most
\[
 348\,\delta_{\mathrm{Line}}(\eps)=348\sqrt\eps.
\]
The answer domain in the second estimate is the completed field alphabet, so
the undefined-evaluation outcome is included.
\end{theorem}
\begin{proof}
\leanok
The point bound is twice the fixed-edge rejection estimate.  For the line
bound, combine the low-degree winning implication with the universal distance
bound $4$ and weaken the resulting linear estimate to the common square-root
error.
\end{proof}
